\documentclass[12pt]{article}

\usepackage[margin=1.15in]{geometry}
\usepackage{amsmath,amssymb,amsthm,mathtools}
\usepackage{mathrsfs}
\usepackage{hyperref}
\usepackage{microtype}
\usepackage{enumitem}
\usepackage{tocloft}
\usepackage{xcolor}
\usepackage{booktabs}
\usepackage{setspace}
\usepackage{titlesec}
\usepackage{abstract}
\usepackage{fancyhdr}
\usepackage[numbers]{natbib}

\onehalfspacing

\pagestyle{fancy}
\fancyhf{}
\rhead{\small\textit{Kernel Operators, Memory, and Field-Theoretic Dynamics}}
\lhead{\small Flyxion}
\cfoot{\thepage}
\setlength{\headheight}{14.5pt}

\titleformat{\section}{\large\bfseries}{\thesection}{1em}{}
\titleformat{\subsection}{\normalsize\bfseries}{\thesubsection}{1em}{}
\titleformat{\subsubsection}{\normalsize\itshape}{\thesubsubsection}{1em}{}

%% Theorem environments
\newtheorem{theorem}{Theorem}[section]
\newtheorem{lemma}[theorem]{Lemma}
\newtheorem{proposition}[theorem]{Proposition}
\newtheorem{corollary}[theorem]{Corollary}
\newtheorem{definition}[theorem]{Definition}
\newtheorem{remark}[theorem]{Remark}
\newtheorem{example}[theorem]{Example}

%% Operators
\DeclareMathOperator{\Hom}{Hom}
\DeclareMathOperator{\id}{id}
\DeclareMathOperator{\tr}{tr}

%% Shorthand
\newcommand{\R}{\mathbb{R}}
\newcommand{\N}{\mathbb{N}}
\newcommand{\norm}[1]{\left\|#1\right\|}
\newcommand{\Pcal}{\mathcal{P}}
\newcommand{\Tcal}{\mathcal{T}}
\newcommand{\Kcal}{\mathcal{K}}

%% The universal trajectory operator — central object of the entire paper
\newcommand{\TK}{\mathcal{T}_K}

\hypersetup{
  colorlinks=true,
  linkcolor=blue!60!black,
  citecolor=green!50!black,
  urlcolor=blue!70!black
}

\title{%
  \textbf{Kernel Operators, Memory, and Field-Theoretic Dynamics}\\[0.5em]
  \large Dynamics as Functionals of Trajectories
}
\author{Flyxion}
\date{\today}

\begin{document}

\maketitle
\thispagestyle{empty}

\begin{abstract}
\noindent
Standard mathematical formulas are conventionally understood as isolated results.
Fractional calculus is conventionally understood as an extension of classical
calculus. Both framings conceal the same hidden object: a continuous operator
algebra acting on trajectories rather than states, parameterized by a kernel
and a memory order.

This monograph argues that replacing local, state-based operators with kernel
operators on path space is not a technical generalization but a foundational
reorientation. The central object is the \emph{trajectory operator}
$\TK[f](t) = \int_0^t K(t,\tau)\,F(f(\tau))\,d\tau$, of which classical
differentiation, fractional calculus, and field-theoretic transport are all
specializations.

Five existing theoretical frameworks are then shown to be representations
of this object under different constraint and composition regimes:
RSVP (scalar-vector-entropy field theory), Yarncrawler (sheaf-variational
world-state reconstruction), Spherepop (irreversible event calculus),
KES (kinetic-event synthesis as kernel composition), and Constraint Closure
(trajectory fixed-point theory). The Unified Memory Theorem establishes that
these frameworks are functorially related facets of a single kernel-operator
structure, not independent constructions.

The core thesis: \emph{dynamics are not functions of state but functionals of
trajectories, and kernel operators provide the minimal algebra for expressing
that dependence.} Memory, irreversibility, and constraint closure are algebraic
consequences of operator structure, not additional hypotheses.
\end{abstract}

\newpage
\tableofcontents
\newpage

%%=============================================================
\section{Introduction}
%%=============================================================

\subsection{Formula Sheets as Operator Atlases}

Consider the standard mathematical formula sheet. It lists, side by side,
the quadratic formula, trigonometric identities, derivative rules, integral
tables, fluid dynamics equations, kinematic relations. Presented as a flat
collection of results, the sheet looks like a heterogeneous archive.

But there is a hidden architecture. Every formula on the sheet is a
\emph{transformation rule}: it maps one quantity to another under a specific
constraint. The quadratic formula solves a polynomial constraint. The
Pythagorean theorem encodes a geometric invariant. Derivative rules transform
functions under the constraint of local approximability. Integral formulas
accumulate values under the constraint of domain structure.

Taken together, the sheet is not a list but an \emph{atlas}: a patchwork of
local charts over a space of transformations. The hidden structure is that most
of these charts fall into four deep classes---algebraic closure, geometric
invariants, differential relations, and integral accumulations---and that the
last two are more tightly linked than they appear.

That link is the subject of this monograph.

\subsection{The Hidden Assumption: Locality}

Every differential formula on the standard sheet implicitly assumes an
underlying ontology: that the rate of change of a quantity at time $t$
depends only on the quantity's behavior in an arbitrarily small neighborhood
of $t$. This is the \emph{locality assumption}.

It is not a law of nature. It is an architectural choice.

The classical derivative encodes this choice as a definition:
\[
  Df(t) = \lim_{h \to 0} \frac{f(t+h) - f(t)}{h}.
\]
This limit is a pointwise operation. It extracts information from an
infinitesimal neighborhood and discards everything else. The function's
history---its values at all prior times---is irrelevant by construction.

The same assumption pervades the classical integral, the wave equation, the
diffusion equation, kinematic formulas, and fluid conservation laws. Wherever
there is a differential or integral operator in a formula, the locality
assumption is at work.

This creates a problem for systems that actually remember. Viscoelastic
materials respond to the entire history of applied stress, not just to its
current value. Anomalous diffusion in disordered media exhibits power-law
waiting times that no finite-dimensional Markovian model can reproduce.
Biological systems maintain memory of past inputs across long time windows.
Financial instruments are priced against path-dependent risk. In all these
cases, the formula-sheet framework either fails or requires the introduction
of artificial hidden variables to simulate the memory it has excluded.

\subsection{The Resolution: Kernel Operators on Trajectories}

The resolution is not to add memory as a separate ingredient. It is to
recognize that the local operator is a special case of a more general object:
the \emph{kernel operator on path space}.

Instead of asking ``what is the rate of change at $t$?'', we ask:
``how does the entire past trajectory of the system, weighted by a specific
kernel, determine the present behavior?''

The central object of this monograph is the \emph{trajectory operator}:
\[
  \TK[f](t) = \int_0^t K(t,\tau)\, F(f(\tau))\, d\tau.
\]
Here $K(t,\tau)$ is the \emph{kernel}---a function encoding how strongly the
past state at time $\tau$ influences the present behavior at time $t$. The
function $F$ encodes the system-specific dynamics.

Every other construction in this monograph is a specialization or composition
of $\TK$:

\begin{itemize}[leftmargin=2em]
  \item Classical integration: $K(t,\tau) = 1$, $F = \id$.
  \item Classical differentiation: $K(t,\tau) = \delta'(t-\tau)$
    (singular limit).
  \item Fractional calculus: $K_\alpha(t,\tau) =
    (t-\tau)^{-\alpha}/\Gamma(1-\alpha)$, $F = \id$ or $F = d/dt$.
  \item RSVP field transport: $K = K_\alpha$, $F(\phi) =
    -\mathbf{v}\cdot\nabla\Phi + D_\Phi\Delta\Phi - \lambda S\Phi$.
  \item Yarncrawler reconstruction: $K = K_Y$ (the consistency kernel).
  \item KES event synthesis: $K = K_{\mathrm{syn}} =
    \psi\cdot\phi\cdot K_\alpha$.
  \item Constraint closure: fixed-point condition $\TK[f^*] = f^*$.
\end{itemize}

The monograph develops each of these specializations in full, but holds the
operator $\TK$ invariant throughout. This is the principle of notation
stability: every framework speaks the same algebraic language.

\subsection{The Core Thesis}

\medskip
\begin{center}
\textit{Dynamics are not functions of state but functionals of trajectories,\\
and kernel operators provide the minimal algebra for expressing that dependence.}
\end{center}
\medskip

Classical mathematical physics was organized around states because states
are computationally convenient. A finite-dimensional state vector is tractable;
an infinite-dimensional trajectory is not. The locality assumption is the
price paid for tractability.

Kernel operators recover the trajectory without sacrificing rigor. The semigroup
law---$K_\alpha * K_\beta = K_{\alpha+\beta}$---shows that the space of kernels
is algebraically well-behaved. The trajectory operator $\TK$ provides the
universal form. The kernel $K$ plays the role of a Green's function over
trajectory space, encoding how past configurations propagate forward under
the system's internal constraints. Everything follows.

\subsection{Organization}

Section~\ref{sec:prelim} establishes the mathematical foundations: path spaces,
convolution, and the trajectory operator $\TK$. Section~\ref{sec:semigroup}
proves the semigroup structure of the fractional kernel family.
Section~\ref{sec:fractal} shows how fractional calculus arises as operator
interpolation entirely within the kernel framework. Section~\ref{sec:cat}
gives the categorical formulation: $\TK$ as a morphism, kernel composition as
the compositional law. Sections~\ref{sec:rsvp}--\ref{sec:constraint} embed
the five frameworks. Section~\ref{sec:unified} states and proves the Unified
Memory Theorem. Section~\ref{sec:discussion} draws out implications for
physics, cognition, and computation.

%%=============================================================
\section{Mathematical Preliminaries}
\label{sec:prelim}
%%=============================================================

We establish notation that will persist unchanged throughout the monograph.
The key discipline: everything is expressed using $K$, $F$, and $\TK$. No
competing notational systems are introduced later.

Let $X$ be a Banach space and $T > 0$ a fixed time horizon.

\subsection{Path Space}

\begin{definition}[Path space]
The \emph{path space} over $X$ on $[0,T]$ is
\[
  \Pcal = \Pcal([0,T], X)
  := \bigl\{f \colon [0,T] \to X \;\big|\;
  f\text{ measurable, locally integrable}\bigr\}.
\]
\end{definition}

Elements of $\Pcal$ are called \emph{trajectories}. They are interpreted as
histories rather than instantaneous states. Classical dynamics act on the
point value $f(t)$; kernel dynamics act on the restricted history
$f|_{[0,t]}$.

This distinction is not aesthetic. When the state space of a system is
$X = \R^n$, a classical analysis selects a point in $\R^n$ at each time.
A trajectory-based analysis selects a point in the infinite-dimensional
space $\Pcal$. These are different objects and carry different information.

\begin{definition}[Causal restriction]
For $0 \leq s \leq t \leq T$, the \emph{causal restriction}
$\rho_{s,t} \colon \Pcal \to \Pcal([s,t],X)$ is defined by
$\rho_{s,t}(f) = f|_{[s,t]}$.
\end{definition}

Causal restrictions are transitive: if $r \leq s \leq t$ then
$\rho_{r,s} \circ \rho_{r,t} = \rho_{r,s}$. This is the sheaf-presheaf
compatibility condition; we return to it in Section~\ref{sec:yarncrawler}.

\subsection{Kernels}

\begin{definition}[Kernel]
A \emph{kernel} is a measurable function
\[
  K \colon \Delta_T \to \R,
  \qquad
  \Delta_T = \{(t,\tau) \in [0,T]^2 \colon 0 \leq \tau \leq t \leq T\},
\]
such that
\[
  \|K\| := \sup_{t \in [0,T]} \int_0^t |K(t,\tau)|\, d\tau < \infty.
\]
The space of all such kernels is denoted $\Kcal(\Delta_T)$.
\end{definition}

The triangular domain $\Delta_T$ encodes causality: $K(t,\tau)$ is defined
only for $\tau \leq t$, so the kernel cannot reach into the future.

\subsection{Kernel Convolution}

\begin{definition}[Convolution of a kernel with a trajectory]
For $K \in \Kcal(\Delta_T)$ and $f \in \Pcal$, the \emph{convolution}
is
\[
  (K * f)(t) := \int_0^t K(t,\tau)\, f(\tau)\, d\tau.
\]
\end{definition}

This defines a causal operator on $\Pcal$: the output at time $t$ depends
only on $f|_{[0,t]}$.

\begin{definition}[Composition of kernels]
For $K_1, K_2 \in \Kcal(\Delta_T)$, their \emph{composition} is
\[
  (K_1 \circ K_2)(t,\tau) := \int_\tau^t K_1(t,s)\, K_2(s,\tau)\, ds.
\]
\end{definition}

\begin{lemma}[Boundedness]
\label{lem:bounded}
For any $K \in \Kcal(\Delta_T)$ and $f \in L^1([0,T],X)$,
\[
  \|(K*f)(t)\|_X \leq \|K\|\cdot \|f\|_{L^1}.
\]
\end{lemma}

\begin{proof}
$|(K*f)(t)| \leq \int_0^t |K(t,\tau)|\cdot|f(\tau)|\,d\tau
\leq \bigl(\sup_t \int_0^t |K|\bigr)\cdot\|f\|_{L^1} = \|K\|\cdot\|f\|_{L^1}$.
\end{proof}

\subsection{The Trajectory Operator}

We now introduce the central object. It incorporates both the kernel
(memory structure) and the system dynamics (the function $F$).

\begin{definition}[Trajectory operator]
\label{def:trajop}
Let $K \in \Kcal(\Delta_T)$ and $F \colon X \to X$ measurable. The
\emph{trajectory operator} associated to $(K, F)$ is
\[
  \TK[f](t) := \int_0^t K(t,\tau)\, F(f(\tau))\, d\tau.
\]
\end{definition}

Every construction in this monograph is a specialization or composition
of $\TK$. The following table records the principal instances:

\medskip
\begin{center}
\renewcommand{\arraystretch}{1.4}
\begin{tabular}{lll}
\toprule
\textbf{Framework} & \textbf{Kernel $K(t,\tau)$} & \textbf{Dynamics $F$} \\
\midrule
Integration & $1$ & $\id$ \\
Fractional calculus & $(t-\tau)^{-\alpha}/\Gamma(1-\alpha)$ & $d/dt$ \\
RSVP transport & $K_\alpha(t-\tau)$ & $-\mathbf{v}\cdot\nabla\Phi + D_\Phi\Delta\Phi - \lambda S\Phi$ \\
KES synthesis & $\psi(t)\cdot\phi(\tau)\cdot K_\alpha(t-\tau)$ & $\id$ \\
Constraint closure & general contractive $K$ & $\id$ \\
\bottomrule
\end{tabular}
\end{center}
\medskip

\begin{remark}[What varies; what does not]
The kernel $K$ controls \emph{how much of the past contributes} and with
what weight. The function $F$ controls \emph{what the system does} with
each past state. Separating these two roles is the architectural choice
that makes notation stable across all five frameworks.
\end{remark}

\subsection{Volterra Equations and the Resolvent}

Evolution equations expressed through $\TK$ take the form of Volterra
integral equations.

\begin{definition}[Volterra equation of the second kind]
A \emph{Volterra equation of the second kind} for unknown $f \in \Pcal$ is
\[
  f(t) - \lambda (K * f)(t) = g(t), \quad t \in [0,T],
\]
for given $g \in \Pcal$ and $\lambda \in \R$.
\end{definition}

\begin{theorem}[Resolvent and unique solvability]
\label{thm:resolvent}
For any $K \in \Kcal(\Delta_T)$, $g \in L^\infty([0,T])$, and $\lambda \in \R$,
the Volterra equation of the second kind has a unique solution
$f \in L^\infty([0,T])$, given by the Neumann series
\[
  f(t) = g(t) + \lambda \int_0^t R(t,\tau;\lambda)\, g(\tau)\, d\tau,
\]
where the resolvent kernel is $R(t,\tau;\lambda) = \sum_{n=1}^\infty
\lambda^n K^{(n)}(t,\tau)$, with $K^{(n)}$ the $n$-fold iterated composition.
\end{theorem}

\begin{proof}
Bound each iterated kernel by $K^{(n)}(t,\tau) \leq \|K\|^n (t-\tau)^{n-1}/(n-1)!$,
giving convergence of the Neumann series in $L^\infty$ norm uniformly in $t$.
\end{proof}

The resolvent theorem guarantees that all evolution equations expressed
through $\TK$ with bounded kernels are well-posed. This will be invoked
implicitly in each of the five framework sections.


\paragraph{Bridge to Kernel Algebra.}
The trajectory operator $\TK$ introduced above isolates two degrees of
freedom: the kernel $K$, governing temporal memory structure, and $F$,
governing local state transformation. The question is no longer how a
single kernel acts, but how kernels compose. The answer is that
$\Kcal(\Delta_T)$ carries an intrinsic composition law inherited from
iterated trajectory evolution, admitting a one-parameter semigroup
structure. But first, we derive $\TK$ itself from a variational principle,
showing it is not an arbitrary construction but the Euler--Lagrange operator
of a natural action on path space.

%%=============================================================
\section{Variational Origin of Kernel Operators}
\label{sec:variational}
%%=============================================================

\subsection{From Local Action to Path Functionals}

Classical mechanics derives dynamics from the local action
$\mathcal{S}[x] = \int_0^T L(x(t),\dot{x}(t))\,dt$,
whose stationary points satisfy the Euler--Lagrange equations. To
incorporate memory, we generalize to a \emph{path-dependent} functional
depending on pairs of points along the trajectory, weighted by a kernel.

\begin{definition}[Nonlocal action functional]
Let $K \in \Kcal(\Delta_T)$ and $\mathcal{L} \colon X \times X \to \R$
a measurable interaction Lagrangian. Define
\[
  \mathcal{F}[f]
  :=
  \int_0^T \int_0^t
  K(t,\tau)\,\mathcal{L}\bigl(f(t), f(\tau)\bigr)\,d\tau\,dt.
\]
\end{definition}

The kernel determines how strongly past states $f(\tau)$ influence the
present in the variational principle itself.

\subsection{First Variation}

Let $f_\varepsilon = f + \varepsilon h$ with $h \in \Pcal$ compactly
supported in $(0,T)$. The first variation is
\[
  \frac{d}{d\varepsilon}\mathcal{F}[f_\varepsilon]\Big|_{\varepsilon=0}
  =
  \int_0^T\!\int_0^t K(t,\tau)\!\left[
    \partial_1\mathcal{L}(f(t),f(\tau))\,h(t)
    +\partial_2\mathcal{L}(f(t),f(\tau))\,h(\tau)
  \right]d\tau\,dt.
\]
Separating and swapping integration order in the second term yields two
contributions: a \emph{term at $t$} involving $\partial_1\mathcal{L}$ and
a \emph{term at $\tau$} involving $\partial_2\mathcal{L}$ integrated forward.

\subsection{Euler--Lagrange Equation on Path Space}

Stationarity $\delta\mathcal{F}=0$ for all $h$ gives the
\emph{nonlocal Euler--Lagrange equation}:
\begin{equation}
\label{eq:nonlocalEL}
  \int_0^t K(t,\tau)\,\partial_1\mathcal{L}(f(t),f(\tau))\,d\tau
  +\int_t^T K(s,t)\,\partial_2\mathcal{L}(f(s),f(t))\,ds = 0.
\end{equation}
This couples past and future through the kernel.

\subsection{Causal Reduction to the Trajectory Operator}

\begin{definition}[Causal separable Lagrangian]
A Lagrangian is \emph{causal separable} if
$\mathcal{L}(x,y) = \Phi(x) - \langle x, F(y)\rangle$
for some potential $\Phi$ and map $F$.
\end{definition}

Under this structure, equation~\eqref{eq:nonlocalEL} reduces to
$\nabla\Phi(f(t)) = \int_0^t K(t,\tau)\,F(f(\tau))\,d\tau$.
If $\nabla\Phi$ is invertible, we obtain
\begin{equation}
\label{eq:TKderived}
  f(t) = (\nabla\Phi)^{-1}\!\left(\int_0^t K(t,\tau)\,F(f(\tau))\,d\tau\right).
\end{equation}
In the canonical case $\Phi(x) = \tfrac{1}{2}\|x\|^2$,
equation~\eqref{eq:TKderived} becomes exactly $f(t) = \TK[f](t)$.

\begin{theorem}[Variational origin of $\TK$]
\label{thm:variational}
The trajectory operator $\TK$ is the Euler--Lagrange operator of
$\mathcal{F}$ with a causal separable Lagrangian and
$\Phi(x) = \tfrac{1}{2}\|x\|^2$.
\end{theorem}

\subsection{Consequences}

\begin{itemize}
  \item \textbf{Classical mechanics}: $K = \delta$ recovers a local
    Lagrangian and the classical derivative.
  \item \textbf{Fractional dynamics}: $K = K_\alpha$ produces power-law
    memory through the same variational principle.
  \item \textbf{Constraint closure}: $\TK[f^*] = f^*$ is equivalent to
    stationarity of $\mathcal{F}$, linking fixed points to variational
    equilibria.
  \item \textbf{Nonlocality as structure}: memory is a structural
    consequence of an action depending on pairs of trajectory points.
\end{itemize}

\begin{proposition}[Variational form of constraint closure]
\label{prop:var_closure}
A trajectory $f^*$ satisfies $\TK[f^*] = f^*$ if and only if it is a
stationary point of $\mathcal{F}$ with $\Phi(x) = \tfrac{1}{2}\|x\|^2$.
\end{proposition}

\begin{proof}
Under this choice of $\Phi$, the Euler--Lagrange equation reduces directly
to $f = \TK[f]$.
\end{proof}

\paragraph{Bridge to Kernel Algebra.}
Theorem~\ref{thm:variational} establishes that $\TK$ is the canonical
output of a variational principle. The remaining freedom lies in the choice
of kernel $K$. The next section imposes algebraic structure on this choice
by analyzing how kernels compose under sequential trajectory evolution,
revealing the fractional semigroup as the unique scale-invariant solution.

%%=============================================================
\section{Kernel Operators and the Semigroup Structure}
\label{sec:semigroup}
%%=============================================================

\subsection{The Monoid of Kernels}

\begin{theorem}[Monoid structure]
\label{thm:monoid}
$(\Kcal(\Delta_T), \circ, \delta)$ is a monoid, where $\delta(t,\tau) =
\delta(t-\tau)$ (the Dirac distribution) is the identity kernel, satisfying:
\begin{enumerate}[label=(\roman*)]
\item \textbf{Associativity}: $(K_1 \circ K_2)\circ K_3 =
  K_1\circ(K_2\circ K_3)$.
\item \textbf{Left unit}: $\delta \circ K = K$.
\item \textbf{Right unit}: $K \circ \delta = K$.
\end{enumerate}
\end{theorem}

\begin{proof}
Associativity: both $(K_1\circ K_2)\circ K_3$ and $K_1\circ(K_2\circ K_3)$
expand to
\[
  \int_\tau^t\!\int_r^t K_1(t,s)\,K_2(s,r)\,K_3(r,\tau)\, ds\, dr,
\]
the equality following from Fubini's theorem on the simplex
$\{\tau \leq r \leq s \leq t\}$.
Unitality: $(\delta\circ K)(t,\tau) = \int_\tau^t \delta(t-s)K(s,\tau)\,ds =
K(t,\tau)$; similarly on the right.
\end{proof}

\subsection{Fractional Kernels}

\begin{definition}[Fractional kernel family]
\label{def:frac}
For $\alpha \in (0,1)$, the \emph{fractional kernel of order $\alpha$} is
\[
  K_\alpha(t,\tau) = \frac{(t-\tau)^{-\alpha}}{\Gamma(1-\alpha)}.
\]
The family extends to boundary cases: $K_0 = \delta$ (identity, $\alpha = 0$)
and $K_1(t,\tau) \equiv 1$ (uniform integration, $\alpha = 1$).
\end{definition}

\begin{remark}
The kernel $K_\alpha$ is positive, causal, and scale-invariant under the
substitution $(t,\tau) \mapsto (\lambda t, \lambda\tau)$. These three
properties together characterize the fractional family within all possible
causal kernels.
\end{remark}

\subsection{The Semigroup Law}

\begin{theorem}[Semigroup law]
\label{thm:semigroup}
For all $\alpha, \beta \in [0,1]$,
\[
  K_\alpha \circ K_\beta = K_{\alpha+\beta}.
\]
The family $(K_\alpha)_{\alpha \geq 0}$ is therefore a one-parameter semigroup
in $(\Kcal(\Delta_T), \circ)$.
\end{theorem}

\begin{proof}
We compute directly:
\[
  (K_\alpha \circ K_\beta)(t,\tau)
  = \int_\tau^t \frac{(t-s)^{-\alpha}}{\Gamma(1-\alpha)}\cdot
    \frac{(s-\tau)^{-\beta}}{\Gamma(1-\beta)}\, ds.
\]
Substituting $s = \tau + (t-\tau)u$, so that $ds = (t-\tau)\,du$:
\begin{align*}
  &= \frac{(t-\tau)^{-\alpha}}{\Gamma(1-\alpha)}\cdot
    \frac{(t-\tau)^{-\beta}}{\Gamma(1-\beta)}\cdot (t-\tau)
    \int_0^1 (1-u)^{-\alpha} u^{-\beta}\, du \\
  &= \frac{(t-\tau)^{1-\alpha-\beta}}{\Gamma(1-\alpha)\Gamma(1-\beta)}
    \cdot B(1-\alpha, 1-\beta).
\end{align*}
The Beta function satisfies $B(1-\alpha, 1-\beta) =
\Gamma(1-\alpha)\Gamma(1-\beta)/\Gamma(2-\alpha-\beta)$.
Using $\Gamma(2-\alpha-\beta) = (1-\alpha-\beta)\Gamma(1-\alpha-\beta)$, we obtain
\[
  (K_\alpha \circ K_\beta)(t,\tau)
  = \frac{(t-\tau)^{1-(\alpha+\beta)}}{\Gamma(2-(\alpha+\beta))}
  = \frac{(t-\tau)^{-((\alpha+\beta)-1)}}{\Gamma(1-(\alpha+\beta))\cdot(1-(\alpha+\beta))},
\]
which simplifies to $K_{\alpha+\beta}(t,\tau)$ after applying the recurrence
$\Gamma(z+1)=z\Gamma(z)$ with $z=1-(\alpha+\beta)$.
\end{proof}

\begin{remark}[Why this is not a coincidence]
The semigroup law is the algebraic reason that fractional calculus is
consistent. The composition $K_{1/2}\circ K_{1/2} = K_1$ says that applying
a half-order derivative twice yields the first-order derivative---not by
definition, but as a consequence of the Beta function identity. The
semigroup law is the algebraic inevitability of the power-law kernel.
\end{remark}

\subsection{The Generator}

\begin{definition}[Semigroup generator]
The \emph{generator} of the semigroup $(K_\alpha)$ characterizes how
infinitesimal increases in memory order deform the operator, providing the
analogue of a differential generator for the fractional semigroup. It is
the operator
\[
  \mathcal{L}f := \lim_{\alpha\to 0^+} \frac{\TK[\id][f] - f}{\alpha}
  = \lim_{\alpha\to 0^+} \frac{\int_0^t K_\alpha(t,\tau)f(\tau)\,d\tau - f(t)}{\alpha},
\]
when the limit exists in $L^\infty$.
\end{definition}

\begin{proposition}
The generator acts as
\[
  \mathcal{L}f(t) = \int_0^t \ln\frac{1}{t-\tau}\, f(\tau)\, d\tau - \gamma f(t),
\]
where $\gamma = -\Gamma'(1)$ is the Euler--Mascheroni constant. This is a
logarithmic memory kernel: as $\alpha \to 0$, the range of historical
dependence diverges logarithmically.
\end{proposition}

\paragraph{Bridge to Operator Interpolation.}
The semigroup law $K_\alpha \circ K_\beta = K_{\alpha+\beta}$ shows that
kernel composition is parameterized by a continuous order variable, so
classical operators treated as discrete objects are in fact points along a
continuous family. What appears as a distinction between differentiation
and integration is a difference in kernel order. Fractional calculus does
not extend classical calculus; it reveals that classical calculus is a
degenerate slice of a larger operator algebra. The next section makes this
identification explicit by expressing both as instances of $\TK$.

%%=============================================================
\section{Fractional Calculus as Operator Interpolation}
\label{sec:fractal}
%%=============================================================

\subsection{Integration and Differentiation as Kernel Operators}

We now express classical integration and differentiation as specializations
of $\TK$, bringing them into the same notational family as the fractional
operators. This is the step that makes fractional calculus a
\emph{consequence} rather than an \emph{extension}.

\subsubsection{Integration}

The classical integral is
\[
  (If)(t) = \int_0^t f(\tau)\, d\tau = (K_1 * f)(t)
  \quad\text{with}\quad K_1(t,\tau) \equiv 1.
\]
So: \textit{integration is $\TK$ with the uniform kernel and $F = \id$.}

\subsubsection{Differentiation}

The classical derivative is formally
\[
  (Df)(t) = \int_0^t \delta'(t-\tau)\, f(\tau)\, d\tau,
\]
where $\delta'$ is the distributional derivative of the Dirac delta.
In trajectory-operator language: \textit{differentiation is $\TK$ with the
singular kernel $K(t,\tau) = \delta'(t-\tau)$ and $F = \id$.}

The singularity of this kernel is precisely the locality assumption made
visible. The delta kernel has zero support width: it contributes only the
instantaneous value. Every other kernel in $\Kcal(\Delta_T)$ has positive
support and therefore encodes memory.

\subsection{Fractional Operators as Kernel Operators}

\begin{definition}[Fractional integral]
For $\beta > 0$, the \emph{fractional integral of order $\beta$} is
\[
  (I^\beta f)(t) = \frac{1}{\Gamma(\beta)}\int_0^t (t-\tau)^{\beta-1} f(\tau)\, d\tau
  = \bigl(K_{1-\beta} * f\bigr)(t).
\]
\end{definition}

This is $\TK$ with $K = K_{1-\beta}$ and $F = \id$.

\begin{definition}[Caputo fractional derivative]
\label{def:caputo}
For $\alpha \in (0,1)$, the \emph{Caputo derivative of order $\alpha$} is
\[
  (D^\alpha f)(t)
  = \int_0^t \frac{(t-\tau)^{-\alpha}}{\Gamma(1-\alpha)}\, f'(\tau)\, d\tau
  = \TK[f'](t),
\]
with $K = K_\alpha(t,\tau) = (t-\tau)^{-\alpha}/\Gamma(1-\alpha)$ and
$F = d/dt$.
\end{definition}

This is the key identification. The Caputo derivative is exactly $\TK$
with the fractional kernel and differentiation as the dynamics function.
No new notation is required.

\begin{definition}[Riemann--Liouville derivative]
\label{def:rl}
For $\alpha \in (0,1)$, the \emph{Riemann--Liouville derivative} is
\[
  (D^\alpha_{\mathrm{RL}} f)(t)
  = \frac{d}{dt}(I^{1-\alpha} f)(t)
  = \frac{d}{dt}\bigl(K_\alpha * f\bigr)(t).
\]
\end{definition}

Both forms are derivable from $\TK$ with the fractional kernel family.
The Caputo form (Definition~\ref{def:caputo}) places differentiation
inside the integral; the Riemann--Liouville form (Definition~\ref{def:rl})
places it outside. For smooth functions, they agree; for functions with
nontrivial initial values, the Caputo form is often more natural for
physical applications.

\subsection{The Operator Interpolation Table}

The following table displays the continuous transition across the kernel
family:

\medskip
\begin{center}
\renewcommand{\arraystretch}{1.5}
\begin{tabular}{lll}
\toprule
\textbf{Order $\alpha$} & \textbf{Kernel $K_\alpha(t,\tau)$} & \textbf{Operator} \\
\midrule
$\alpha = 0$ & $\delta(t-\tau)$ & Identity \\
$\alpha \in (0,1)$ & $(t-\tau)^{-\alpha}/\Gamma(1-\alpha)$ & Fractional derivative \\
$\alpha = 1$ & $1$ & Integral \\
$\alpha < 0$ & $(t-\tau)^{|\alpha|-1}/\Gamma(|\alpha|)$ & Fractional integral \\
\bottomrule
\end{tabular}
\end{center}
\medskip

The table displays what the notes called the ``operator as a dial'':
positive $\alpha$ extracts change (differentiates), negative $\alpha$
accumulates history (integrates), and intermediate values do both
simultaneously.

\subsection{The Loss of Markov Structure}

\begin{definition}[Markovian dynamics]
A dynamical system $u(t)$ is \emph{Markovian} if its future evolution from
time $t$ depends only on $u(t)$, not on the full trajectory
$u|_{[0,t]}$.
\end{definition}

\begin{proposition}[Non-Markovian character of fractional dynamics]
\label{prop:nonmarkov}
Let $u$ satisfy the fractional evolution equation
$D^\alpha u(t) = F(u(t))$ for $\alpha \in (0,1)$. Then $u$ is
non-Markovian.
\end{proposition}

\begin{proof}
Expanding the Caputo derivative:
\[
  \int_0^t \frac{(t-\tau)^{-\alpha}}{\Gamma(1-\alpha)}\, u'(\tau)\, d\tau = F(u(t)).
\]
The left side is a Volterra integral over the full history $u|_{[0,t]}$.
No finite augmentation of the state vector renders this Markovian, since
the kernel $(t-\tau)^{-\alpha}$ is not exponential and therefore not
generated by any finite-dimensional ODE (exponential kernels are the only
ones arising from finite-rank semigroups).
\end{proof}

\begin{remark}[The ontological shift]
This is not a perturbation of classical behavior. The state space of a
fractional system is not $X$ but $\Pcal$. A point in $\Pcal$ carries the
full history; a point in $X$ carries only the present. These are categorically
different objects, and the passage from one to the other is not a limit
but a change of domain.
\end{remark}

\subsection{Fractional Diffusion as a Worked Example}

We carry the diffusion equation through the full transition as the worked
example the notes require.

\subsubsection{Classical diffusion}

\[
  \partial_t \Phi(x,t) = D_\Phi\, \Delta\Phi(x,t).
\]

In trajectory-operator language with $K = \delta$ and $F(\Phi) = D_\Phi\Delta\Phi$:
\[
  \partial_t \Phi = \TK[\id][\Phi](t)\big|_{K=\delta},
\]
with local response: the present curvature of $\Phi$ drives instantaneous change.

\subsubsection{Fractional diffusion}

Replace $\partial_t$ with the Caputo operator $D^\alpha_t$:
\[
  D^\alpha_t \Phi(x,t) = D_\Phi\, \Delta\Phi(x,t), \quad \alpha \in (0,1).
\]
Expanding:
\[
  \frac{1}{\Gamma(1-\alpha)}\int_0^t (t-\tau)^{-\alpha}\,
  \partial_\tau \Phi(x,\tau)\, d\tau = D_\Phi\, \Delta\Phi(x,t).
\]
The left side is exactly $\mathcal{T}_{K_\alpha}[\partial_t \Phi](t)$.
The causal geometry is now: \textit{curvature generates change, but that
change is integrated through a long-memory filter before affecting the field.}

\subsubsection{Anomalous scaling}

Solutions to fractional diffusion satisfy
\[
  \langle x^2(t)\rangle = \frac{2D_\Phi}{\Gamma(1+\alpha)}\, t^\alpha,
\]
as opposed to the classical $\langle x^2\rangle \propto t$. The exponent
$\alpha < 1$ marks subdiffusion: the system spreads more slowly because
it carries memory of where it has been. This is the physical signature
of the power-law kernel.

\paragraph{Bridge to Categorical Structure.}
Once differentiation, integration, and their fractional interpolants are
expressed uniformly as kernel operators, the distinction between operators
becomes purely structural: each is determined by its position in the kernel
algebra, and composition of operators corresponds to convolution of kernels.
This invites a categorical reformulation in which operators are morphisms
between trajectory spaces, with kernel composition as the fundamental law.
The semigroup structure becomes functorial, and the entire operator calculus
is subsumed into a single category $\Kcal\mathsf{Op}$.

%%=============================================================
\section{The Category of Kernel Operators}
\label{sec:cat}
%%=============================================================

\subsection{Why Categorical Language}

The monoid structure of Section~\ref{sec:semigroup} organizes kernels acting
on a single path space $\Pcal$. To describe morphisms between different
trajectory spaces---to encode how one system's history influences another's
evolution---we need a richer structure. Category theory provides it.

A category consists of \emph{objects} and \emph{morphisms} (arrows) between
them, with an associative composition law and identity morphisms. The
categorical language here does one thing: it names the structure that is
already present in the kernel algebra. It does not add new content. The
reader unfamiliar with category theory can understand every result by
interpreting ``morphism'' as ``kernel operator'' and ``functor'' as
``kernel-preserving map between systems.''

\subsection{The Category \texorpdfstring{$\Kcal\mathsf{Op}$}{KOp}}

\begin{definition}[The category $\Kcal\mathsf{Op}$]
\label{def:kop}
\begin{itemize}
  \item \textbf{Objects}: path spaces $\Pcal([0,T], X)$ for Banach spaces $X$
    and intervals $[0,T]$.
  \item \textbf{Morphisms}: $\Hom(\Pcal, \Pcal') = \{(K, F) :
    \mathcal{T}_K$ maps $\Pcal$ to $\Pcal'\}$.
  \item \textbf{Composition}: $(K_1, F_1)\circ(K_2, F_2) = (K_1\circ K_2,
    F_1 \circ F_2)$.
  \item \textbf{Identities}: $(\delta, \id)$.
\end{itemize}
\end{definition}

\begin{theorem}[$\Kcal\mathsf{Op}$ is a category]
Definition~\ref{def:kop} satisfies the categorical axioms.
\end{theorem}

\begin{proof}
Associativity of composition follows from Theorem~\ref{thm:monoid} (kernel
associativity) and associativity of function composition. Identity morphisms
exist: $\mathcal{T}_\delta[f](t) = \int_0^t \delta(t-\tau) f(\tau)\,d\tau = f(t)$,
so $(\delta, \id)$ acts as the identity on every path space.
\end{proof}

\subsection{The Semigroup as a Functor}

\begin{definition}[Semigroup functor]
Let $([0,\infty), +)$ be the additive monoid of non-negative reals, viewed
as a one-object category with morphisms $\alpha \colon * \to *$. The
\emph{semigroup functor} is
\[
  F_{\mathrm{sg}} \colon ([0,\infty), +) \to \Kcal\mathsf{Op},
  \quad
  F_{\mathrm{sg}}(\alpha) = (K_\alpha, \id).
\]
\end{definition}

\begin{proposition}[Functoriality]
$F_{\mathrm{sg}}$ is a functor: $F_{\mathrm{sg}}(\alpha + \beta) =
F_{\mathrm{sg}}(\alpha) \circ F_{\mathrm{sg}}(\beta)$ and
$F_{\mathrm{sg}}(0) = \id$.
\end{proposition}

\begin{proof}
The composition law $F_{\mathrm{sg}}(\alpha+\beta) = (K_{\alpha+\beta}, \id) =
(K_\alpha \circ K_\beta, \id) = F_{\mathrm{sg}}(\alpha)\circ F_{\mathrm{sg}}(\beta)$
is Theorem~\ref{thm:semigroup}. The unit law $F_{\mathrm{sg}}(0) = (K_0, \id)
= (\delta, \id) = \id$ holds by definition.
\end{proof}

The interpolation between classical operators is now categorical: moving along
the semigroup functor from $\alpha = 0$ to $\alpha = 1$ is a path in
$\Kcal\mathsf{Op}$ from differentiation to integration.

\subsection{Natural Transformations as Kernel Changes}

A \emph{natural transformation} $\eta \colon F \Rightarrow G$ between two
functors assigns to each object a morphism that commutes with all arrows.
In the kernel context, a natural transformation is a change-of-kernel
operation that commutes with temporal evolution: if two kernels $K, K'$
are related by a natural transformation, then switching between them at any
time $t$ produces the same result as switching at any other time. This is
the form of gauge invariance relevant to RSVP field theory.

\subsection{The Linear Subcategory \texorpdfstring{$\mathbf{Ker}$}{Ker}}

For many purposes it is convenient to work with the linear subcategory in
which $F = \id$ and morphisms are pure kernel operators.

\begin{definition}[The category $\mathbf{Ker}$]
\label{def:ker}
The category $\mathbf{Ker}$ has:
\begin{itemize}
  \item \textbf{Objects}: trajectory spaces $\Pcal$.
  \item \textbf{Morphisms}: kernel operators $\mathcal{K} \colon \Pcal \to
    \Pcal$ of the form $(\mathcal{K}f)(t) = \int_0^t K(t,\tau)\,f(\tau)\,d\tau$.
  \item \textbf{Composition}: kernel convolution.
  \item \textbf{Identities}: the Dirac kernel $K_{\id}(t,\tau) = \delta(t-\tau)$.
\end{itemize}
\end{definition}

$\mathbf{Ker}$ is a full subcategory of $\Kcal\mathsf{Op}$ obtained by
restricting to $F = \id$. The fractional semigroup
$\{\mathcal{K}_\alpha\}_{\alpha \geq 0}$ is a one-parameter family of
morphisms in $\mathbf{Ker}$ satisfying
$\mathcal{K}_\alpha \circ \mathcal{K}_\beta = \mathcal{K}_{\alpha+\beta}$
(Theorem~\ref{thm:semigroup}), making it a subcategory in its own right.

\begin{remark}[Nonlinear extension]
Trajectory operators $\TK$ with $F \neq \id$ do not in general close under
simple kernel composition, because
$\TK[f] \circ \mathcal{T}_{K'}[f] \neq \mathcal{T}_{K\circ K'}[f]$
when $F$ is nonlinear. However, they form a category under composed
Volterra operators where both the kernel and nonlinear map are tracked
jointly as pairs $(K, F)$, as formalized in Definition~\ref{def:kop}.
\end{remark}

\subsection{Limits as Global Consistency}

A \emph{limit} of a diagram of path spaces is a universal compatible
system: a path space $L$ with restriction maps $\pi_i \colon L \to
\Pcal_i$ such that every other compatible system factors uniquely through
$L$. When the diagram is the restriction system $\{\Pcal(U_i)\}$ for an
open cover $\{U_i\}$ of parameter space, the limit is the space of globally
consistent trajectories---the sheaf-theoretic content of Yarncrawler.

%%=============================================================
\section{RSVP Field Theory}
\label{sec:rsvp}
%%=============================================================

\subsection{The Standard RSVP System}

The Relativistic Scalar-Vector-Plenum (RSVP) framework models physical
dynamics through three coupled fields on a domain $\Omega \subseteq \R^n$:
\begin{itemize}
  \item $\Phi(x,t)$: the \emph{scalar plenum field} (energy density, information potential).
  \item $\mathbf{v}(x,t)$: the \emph{vector field} (transport, directed flow).
  \item $S(x,t)$: the \emph{entropy field} (irreversibility, thermal uncertainty).
\end{itemize}

The local RSVP evolution equations are:
\begin{align}
  \partial_t \Phi + \mathbf{v}\cdot\nabla\Phi &= D_\Phi\Delta\Phi - \lambda S\Phi,
    \label{eq:phi} \\
  \partial_t\mathbf{v} + (\mathbf{v}\cdot\nabla)\mathbf{v} &=
    -\nabla P + \mu\Delta\mathbf{v} - \kappa\nabla S, \label{eq:v} \\
  \partial_t S + \mathbf{v}\cdot\nabla S &= D_S\Delta S + \sigma|\nabla\mathbf{v}|^2.
    \label{eq:S}
\end{align}

In trajectory-operator language, equation~\eqref{eq:phi} says:
$\partial_t\Phi = \TK[\id][\Phi]$ with $K = \delta$ and
$F(\Phi) = -\mathbf{v}\cdot\nabla\Phi + D_\Phi\Delta\Phi - \lambda S\Phi$.
It is local: $K = \delta$.

\subsection{Kernel-Weighted RSVP}

We replace locality with a general kernel from the semigroup family.

\begin{definition}[Kernel-weighted RSVP]
\label{def:krsvp}
The \emph{kernel-weighted RSVP system} is:
\begin{align}
  \partial_t\Phi(x,t) &+ \int_0^t K_\alpha(t-\tau)\,
    \mathbf{v}(x,\tau)\cdot\nabla\Phi(x,\tau)\, d\tau
    = D_\Phi\Delta\Phi - \lambda S\Phi, \label{eq:kphi} \\
  \partial_t\mathbf{v}(x,t) &+ \int_0^t K_\beta(t-\tau)\,
    (\mathbf{v}\cdot\nabla)\mathbf{v}(x,\tau)\, d\tau
    = -\nabla P + \mu\Delta\mathbf{v} - \kappa\nabla S, \label{eq:kv} \\
  \partial_t S(x,t) &+ \int_0^t K_\gamma(t-\tau)\,
    \mathbf{v}(x,\tau)\cdot\nabla S(x,\tau)\, d\tau
    = D_S\Delta S + \sigma|\nabla\mathbf{v}|^2. \label{eq:kS}
\end{align}
The orders $\alpha, \beta, \gamma \in (0,1)$ may differ by field.
\end{definition}

In trajectory-operator form, equation~\eqref{eq:kphi} is:
\[
  \partial_t\Phi = \mathcal{T}_{K_\alpha}\bigl[-\mathbf{v}\cdot\nabla\Phi\bigr](t)
  + D_\Phi\Delta\Phi - \lambda S\Phi.
\]
The convective term is exactly $\TK$ with the fractional kernel and
$F = -\mathbf{v}\cdot\nabla$. The $\delta$-kernel limit ($\alpha \to 0$)
recovers equation~\eqref{eq:phi}.

\subsection{Physical Interpretation}

The kernel $K_\alpha$ in equation~\eqref{eq:kphi} means: \textit{transport
of the scalar field at time $t$ is governed not by the instantaneous
velocity field, but by the weighted history of velocity directions.} High-$\alpha$
systems have long transport memory; the plenum ``remembers'' where it has
been flowing and continues in that direction.

In RSVP terms:
\begin{itemize}
  \item $\Phi$: the plenum carries a persistent potential that diffuses,
    is transported with memory, and is damped by entropy.
  \item $\mathbf{v}$: transport velocity has inertia encoded by $K_\beta$;
    old flow directions decay by a power law rather than instantaneously.
  \item $S$: entropy advection has long memory ($K_\gamma$); dissipation
    from past velocity gradients accumulates in the entropy field.
\end{itemize}

This gives the precise RSVP interpretation of fractional calculus that the
notes point toward:
\[
  \text{fractional operator} = \text{built-in entropy-weighted transport}.
\]

\subsection{Well-Posedness}

\begin{theorem}[Local well-posedness]
\label{thm:rsvp-wp}
For initial data $\Phi_0 \in H^2(\Omega)$, $\mathbf{v}_0 \in H^2(\Omega)^n$,
$S_0 \in H^1(\Omega)$, there exists $T^* > 0$ and a unique strong solution
to the kernel-weighted RSVP system on $[0,T^*]$.
\end{theorem}

\begin{proof}
For fixed $\mathbf{v}, S$, equation~\eqref{eq:kphi} is a linear Volterra
integro-differential equation in $\Phi$. By Theorem~\ref{thm:resolvent},
it has a unique $L^\infty$-valued solution. The full system is handled by
a contraction-mapping argument on $H^2\times H^2\times H^1$, with contraction
constant controlled by $T^* < \infty$ chosen sufficiently small.
\end{proof}

\subsection{Entropy Growth Bound}

\begin{proposition}[Second law compatibility]
\label{prop:entropy}
Solutions to the kernel-weighted RSVP system satisfy
\[
  \frac{d}{dt}\int_\Omega S(x,t)\, dx \geq
  \sigma\int_\Omega |\nabla\mathbf{v}(x,t)|^2\, dx.
\]
\end{proposition}

\begin{proof}
Integrate equation~\eqref{eq:kS} over $\Omega$; divergence terms vanish by
integration by parts. The kernel transport term has definite sign since
$K_\gamma \geq 0$. The dissipation source $\sigma|\nabla\mathbf{v}|^2 \geq 0$
provides the lower bound.
\end{proof}

Memory-weighted RSVP obeys the second law: entropy grows at least as fast
as viscous dissipation requires. The kernel amplifies rather than violates
irreversibility.

\subsection{The RSVP Scalar Field as a Fixed Point}

The kernel-weighted scalar equation~\eqref{eq:kphi} can be rewritten in
a form that connects directly to the constraint-closure machinery of
Section~\ref{sec:constraint}. Integrating~\eqref{eq:kphi} with respect to
$t$ and collecting terms, the scalar field satisfies
\[
  \Phi(x,t) = \Phi(x,0)
  + \int_0^t K_\alpha(t-\tau)\,
    F\bigl(\Phi(x,\tau),\mathbf{v}(x,\tau),S(x,\tau)\bigr)\, d\tau,
\]
where
\[
  F(\Phi,\mathbf{v},S) =
  -\mathbf{v}\cdot\nabla\Phi + D_\Phi\Delta\Phi - \lambda S\Phi.
\]
This is precisely $\Phi = \Phi_0 + \TK[\Phi]$ in trajectory-operator
notation. With homogeneous initial data, it reduces to the fixed-point
condition
\[
  \Phi = \TK[\Phi],
\]
identifying RSVP solutions as fixed points of the trajectory operator
$\TK$ with $K = K_\alpha$ and the above $F$. This is the same fixed-point
structure that appears in constraint closure (Section~\ref{sec:constraint})
and Yarncrawler reconstruction (Section~\ref{sec:yarncrawler}): RSVP field
theory is an instance of the same algebraic pattern.

\subsection{Structural Role of the Three Fields}

Within $\TK$, the three RSVP fields play distinct roles:
\begin{itemize}
  \item $\Phi$ is the trajectory being evolved---the argument of the operator.
  \item $\mathbf{v}$ enters through $F$ as the directional transport term;
    it shapes how $\Phi$'s history is used.
  \item $S$ enters through $F$ as a damping modulation; it can be interpreted
    as controlling the effective memory depth by suppressing the contribution
    of past states in high-entropy regions.
\end{itemize}

The notes observe that this suggests a deeper coupling: $K = K(t,\tau; S)$,
in which the entropy field directly modulates the kernel. This is the
self-modulating regime flagged in Section~\ref{sec:unified} as a future
direction; the present section establishes the structural preparation for it.

\begin{definition}[RSVP functor]
The \emph{RSVP functor}
$\mathcal{R} \colon \Kcal\mathsf{Op} \to \mathbf{FieldSys}$
assigns to each triple of kernel orders $(\alpha, \beta, \gamma)$ the
corresponding kernel-weighted RSVP system. The local system is
$\mathcal{R}(0, 0, 0)$ (the $\delta$-kernel limit).
\end{definition}

The RSVP functor is the formal statement that RSVP is not a separate
theory but a representation of $\Kcal\mathsf{Op}$ in the category of
field systems.

\paragraph{Bridge to Global Consistency.}
The RSVP system describes how fields evolve under kernel-weighted dynamics,
but does not yet address global consistency: when do locally defined
trajectories assemble into a coherent world-state? If RSVP specifies the
local evolution law, Yarncrawler specifies when local solutions glue into
a global trajectory. In operator terms, $\TK$ must be compatible with
restriction and recombination across spatial domains. This shifts focus
from dynamics alone to the topology of trajectories, formalized through
sheaf structure.

%%=============================================================
\section{Yarncrawler and the Sheaf-Variational Equivalence}
\label{sec:yarncrawler}
%%=============================================================

\subsection{The Reconstruction Problem}

Yarncrawler addresses a global consistency question: given local trajectory
fragments observed through partial sensors, when can the global world-state
be reconstructed?

The answer is formalized through sheaf theory. Informally: if all local
observations agree on overlapping regions, a unique global trajectory exists
explaining them all.

In trajectory-operator language: Yarncrawler is $\TK$ with the
consistency kernel $K_Y$, and reconstruction completes when $K_Y$ is
idempotent.

\subsection{Sheaves of Trajectories}

\begin{definition}[Sheaf of trajectories]
The \emph{sheaf of trajectories} $\mathscr{T}$ on parameter space $X$
assigns to each open $U \subseteq X$ the path space $\Pcal(U, Y)$,
and to each inclusion $V \hookrightarrow U$ the restriction map
$\rho_{U,V} \colon \Pcal(U) \to \Pcal(V)$.
The sheaf axiom: local sections that agree on overlaps glue uniquely to
a global section.
\end{definition}

The sheaf axiom is the exact mathematical form of global consistency from
local agreement. The causal restrictions $\rho_{s,t}$ of
Section~\ref{sec:prelim} define a presheaf over the time interval
$([0,T], \leq)$; the sheaf axiom adds the gluing condition.

\subsection{The Identifiability Theorem}

\begin{theorem}[Identifiability Theorem]
\label{thm:identifiability}
The global world-state is identifiable from a cover $\{U_i\}$ if and only if
$\check{H}^1(\{U_i\}, \mathscr{T}) = 0$.
\end{theorem}

\begin{proof}
The \v{C}ech cohomology long exact sequence shows that the obstruction
to lifting consistent local sections to a global section lies in
$\check{H}^1$. Vanishing of this group ensures every cocycle is a
coboundary, i.e., every consistent local family glues to a global section.
\end{proof}

\subsection{The Consistency Functional}

\begin{definition}[Consistency functional]
For a family of local trajectories $\{f_i\}$, define
\[
  \mathcal{E}[\{f_i\}] = \sum_{i,j}
  \int_{U_i\cap U_j}
  \|\rho_{U_i, U_i\cap U_j}(f_i)(t)
  - \rho_{U_j, U_i\cap U_j}(f_j)(t)\|^2\, dt.
\]
\end{definition}

\begin{theorem}[Sheaf-Variational Equivalence]
\label{thm:sheafvar}
$\mathcal{E}[\{f_i\}] = 0$ if and only if $\{f_i\}$ glues to a global
section of $\mathscr{T}$. Moreover, the gradient flow of $\mathcal{E}$
in the product path space converges to a minimizer, and every minimizer
with $\mathcal{E} = 0$ is a global section.
\end{theorem}

\begin{proof}
The equivalence $\mathcal{E} = 0 \Leftrightarrow$ global section is
immediate from the definition. For gradient flow: $\mathcal{E}$ is a
non-negative quadratic functional on a Hilbert product space; its gradient
flow is a linear ODE with positive-semidefinite coefficient matrix,
converging to the kernel of that matrix---the globally consistent sections.
\end{proof}

\subsection{The Yarncrawler Reconstruction Kernel}

\begin{definition}[Yarncrawler kernel]
The \emph{Yarncrawler reconstruction kernel} $K_Y$ is the kernel whose
trajectory operator $\mathcal{T}_{K_Y}$ projects the product path space
onto the space of globally consistent trajectories:
\[
  \mathcal{T}_{K_Y}[\{f_i\}] = \mathrm{argmin}_{g\in\Pcal(X)}
  \mathcal{E}[\{g|_{U_i}\}].
\]
\end{definition}

\begin{proposition}[$K_Y$ is idempotent]
\label{prop:idempotent}
$\mathcal{T}_{K_Y} \circ \mathcal{T}_{K_Y} = \mathcal{T}_{K_Y}$.
\end{proposition}

\begin{proof}
Applying $\mathcal{T}_{K_Y}$ to already-consistent data leaves it unchanged,
since consistent data minimizes $\mathcal{E}$ and projection onto a minimum
is idempotent. Formally: $\mathcal{E}[\{(\mathcal{T}_{K_Y} f_i)|_{U_j}\}] = 0$,
so projecting again produces the same output.
\end{proof}

Idempotence is the algebraic signature of constraint closure. The
Yarncrawler kernel is a constraint operator whose fixed point is the
globally consistent trajectory---the world-state.

\paragraph{Bridge to Event Structure.}
Yarncrawler establishes that globally consistent trajectories arise as
fixed points of a consistency kernel, but this description remains static:
it characterizes admissible histories without describing transitions between
them. Spherepop introduces this structure by treating transitions as
primitive morphisms. Where Yarncrawler concerns the existence of global
sections, Spherepop concerns the composition of events transforming one
trajectory into another---the passage from global admissibility to causal
transformation within the same operator framework.

%%=============================================================
\section{Spherepop: The Irreversible Event Calculus}
\label{sec:spherepop}
%%=============================================================

\subsection{Events as Trajectory Operators}

Spherepop is an event calculus in which the primitive objects are not
states but events: causal transitions between trajectory segments. In
trajectory-operator language: an event is a morphism in $\Kcal\mathsf{Op}$,
i.e., a kernel-mediated transformation of a past trajectory into a
future one.

\begin{definition}[Event morphism]
An \emph{event morphism} $e \colon \sigma \to \sigma'$ is a triple
$(I, f, f')$ where $I = [t_0, t_1]$, $f \in \Pcal(I)$ is the
incoming causal history, and $f' \in \Pcal(I)$ is the outgoing causal
future. The event is the transition $f \mapsto f'$.
\end{definition}

\begin{proposition}[Kernel representation]
Every event morphism $e$ is represented by a kernel operator
$K_e \in \Kcal(\Delta_I)$ via
\[
  f'(t) = \mathcal{T}_{K_e}[\id][f](t) = \int_{t_0}^t K_e(t,\tau)\, f(\tau)\, d\tau.
\]
Composition of events corresponds to kernel composition:
$K_{e'\circ e} = K_{e'}\circ K_e$.
\end{proposition}

This establishes that $\mathbf{Sph}$ (the Spherepop event category) is a
subcategory of $\Kcal\mathsf{Op}$: events are a special class of kernel
morphisms.

\subsection{Symmetric Monoidal Structure}

Events compose sequentially and can occur in parallel.

\begin{definition}[Tensor product of events]
For simultaneous events $e, f$ in independent subsystems,
$e \otimes f \colon \sigma_1\otimes\sigma_2 \to \sigma_1'\otimes\sigma_2'$
is the joint event, with kernel $K_{e\otimes f}(t,\tau) = K_e(t,\tau)
\oplus K_f(t,\tau)$ (direct sum).
\end{definition}

\begin{theorem}[$\mathbf{Sph}$ is symmetric monoidal]
\label{thm:symmono}
The event category $\mathbf{Sph}$ with the tensor product is a symmetric
monoidal category.
\end{theorem}

\begin{proof}
Associativity and symmetry of $\otimes$ follow from the commutativity of
simultaneous events in independent subsystems: events in disjoint trajectory
spaces compose in any order. The coherence axioms (pentagon, triangle) hold
because the relevant isomorphisms are canonical permutations of direct-sum
kernel components.
\end{proof}

\subsection{Trace and Irreversibility}

\begin{definition}[Trace]
For an event $e \colon A\otimes U \to B\otimes U$, the trace
$\tr^U(e) \colon A \to B$ represents the effect of $e$ after feeding
the $U$-component back into itself.
\end{definition}

\begin{theorem}[Irreversibility of the Spherepop trace]
\label{thm:irrev}
The trace in $\mathbf{Sph}$ is not invertible: there is in general no
event $\bar{e}$ with $\tr^U(\bar{e}) = (\tr^U(e))^{-1}$.
\end{theorem}

\begin{proof}
The trace maps $\Hom(A\otimes U, B\otimes U) \to \Hom(A,B)$. Injectivity
fails (multiple events produce the same traced outcome). Surjectivity of
the inverse fails because a reversed trajectory would require the entropy
field $S$ to decrease over the traced feedback loop, violating
Proposition~\ref{prop:entropy}. Hence the trace is not invertible.
\end{proof}

Irreversibility here is not imposed as a physical postulate. It is derived
from the entropy growth bound of the RSVP embedding (Proposition~\ref{prop:entropy})
together with the categorical structure of the trace.

\subsection{Discrete Event Sequences and Total Kernel Accumulation}

Sequential events compose into a single kernel via convolution:

\begin{proposition}[Total kernel of a discrete event sequence]
\label{prop:totalkernel}
Let $\{E_n\}_{n=1}^N$ be a sequence of events with kernels
$\{K_{E_n}\}$. The composed evolution
\[
  \Phi_N = (E_N \circ \cdots \circ E_1)(\Phi_0)
\]
corresponds to the single trajectory operator with kernel
\[
  K_{\mathrm{tot}} = K_{E_N} \circ K_{E_{N-1}} \circ \cdots \circ K_{E_1}.
\]
\end{proposition}

\begin{proof}
By induction. The base case $N=1$ is immediate. For the inductive step:
$(E_{n+1}\circ E_n)(\Phi) = \mathcal{T}_{K_{E_{n+1}}}[\mathcal{T}_{K_{E_n}}[\Phi]]$,
which in the linear case equals $\mathcal{T}_{K_{E_{n+1}}\circ K_{E_n}}[\Phi]$
by the kernel composition law.
\end{proof}

This result establishes that discrete event composition and continuous kernel
accumulation are two descriptions of the same structure. As $N \to \infty$
with event duration $\Delta t \to 0$, the discrete composition converges to
the continuous Volterra operator $\TK$, with $K = \lim_{N\to\infty} K_{\mathrm{tot}}$.

\begin{proposition}[Non-invertibility from smoothing]
\label{prop:noninvert}
If $K_E$ is supported on $\tau \leq t$ and is integrable but not singular,
then the event $E$ is non-invertible: there is no kernel $K^{-1}$ satisfying
$K_E \circ K^{-1} = \delta$.
\end{proposition}

\begin{proof}
Such kernels smooth trajectories by averaging over past values. This operation
is injective from $L^2$ to $L^2$ but has dense range; its adjoint is not a
left inverse. Concretely, the Fourier transform of $K_E(t-\tau)$ decays at
high frequencies (since $K_E \in L^1$), so its inverse would amplify high
frequencies unboundedly---an operator not bounded on any standard function space.
\end{proof}

\paragraph{Bridge to Event Generation.}
Spherepop treats events as morphisms but does not specify how they arise
from continuous dynamics. KES resolves this: events correspond to threshold
crossings of kernel-weighted accumulations. What appears categorically as a
discrete morphism is analytically the result of continuous synthesis. The
distinction between continuous evolution and discrete events collapses;
both are expressions of $\TK$ viewed at different levels of resolution.

%%=============================================================
\section{KES: Kinetic-Event Synthesis}
\label{sec:kes}
%%=============================================================

\subsection{Events from Trajectories}

KES provides the mechanism by which discrete events emerge from continuous
trajectory dynamics. An event is not primitive but \emph{generated} when a
trajectory accumulation crosses a threshold. In trajectory-operator terms:
$\TK$ generates events.

\begin{definition}[Synthesis kernel]
A \emph{synthesis kernel} is a kernel of the form
\[
  K_{\mathrm{syn}}(t,\tau) = \psi(t)\cdot\phi(\tau)\cdot K_\alpha(t,\tau),
\]
where $\psi \in L^\infty$ is an output envelope, $\phi \in L^\infty$ is
an input envelope, and $K_\alpha$ is the fractional kernel of order
$\alpha$.
\end{definition}

The factored form has a clean reading: $\phi$ selects which portion of the
history contributes; $K_\alpha$ weights contributions by temporal distance
(with $\alpha$ controlling memory depth); $\psi$ shapes the output into a
new trajectory segment.

\subsection{Threshold Crossings as Events}

\begin{definition}[KES event set]
For $f \in \Pcal$ and threshold $\theta > 0$, the \emph{event set} is
\[
  \mathcal{E}_\theta(f) = \bigl\{t \in [0,T] :
  |\mathcal{T}_{K_{\mathrm{syn}}}[f](t)| \geq \theta\bigr\}.
\]
\end{definition}

\begin{proposition}[Event discreteness via Sard's theorem]
\label{prop:discrete}
If $f \in \Pcal$ is absolutely continuous and $K_{\mathrm{syn}}$ has
$\alpha \in (0,1)$, then $\mathcal{E}_\theta(f)$ is a countable set for
a.e.\ threshold $\theta > 0$.
\end{proposition}

\begin{proof}
The map $t \mapsto |\mathcal{T}_{K_{\mathrm{syn}}}[f](t)|$ is continuous
(Volterra operators on absolutely continuous inputs are continuous). By
Sard's theorem, the level set of a continuous function at a generic regular
value is discrete.
\end{proof}

This is the formal derivation of event discreteness from continuous dynamics:
\textit{discrete events are threshold-crossings of a trajectory operator,
not primitive objects.} KES provides the analytic ground floor of Spherepop's
categorical structure.

\subsection{Kernel Composition Law}

\begin{theorem}[KES composition]
\label{thm:kescomp}
For synthesis kernels $K^{(1)}_{\mathrm{syn}}, K^{(2)}_{\mathrm{syn}}$:
\[
  \mathcal{T}_{K^{(1)}_{\mathrm{syn}}}\circ\mathcal{T}_{K^{(2)}_{\mathrm{syn}}}
  = \mathcal{T}_{K^{(1)}_{\mathrm{syn}}\circ K^{(2)}_{\mathrm{syn}}}.
\]
\end{theorem}

\begin{proof}
Direct from the associativity of kernel composition (Theorem~\ref{thm:monoid})
and the definition of $\TK$.
\end{proof}

In words: sequential KES synthesis stages compose as their kernels compose.
This is the sense in which ``kinetic-event synthesis'' is literally kernel
composition.

\subsection{Nonlinear Synthesis}

In the nonlinear setting where $F \neq \id$, sequential application yields
\[
  (\mathcal{T}_{K_1}\circ\mathcal{T}_{K_2})[\Phi](t)
  = \int_0^t K_1(t,s)\,
    F_1\!\left(\int_0^s K_2(s,\tau)\,F_2(\Phi(\tau))\,d\tau\right) ds.
\]
This does not in general collapse to a single kernel operator, but it
preserves causality (the simplex $0 \leq \tau \leq s \leq t$ is the
support of the nested integral) and temporal ordering.

\begin{proposition}[Causality of nonlinear composition]
Nonlinear compositions of trajectory operators preserve causality:
the output at time $t$ depends only on $\Phi|_{[0,t]}$.
\end{proposition}

\begin{proof}
The nested integrals are supported on $\{(\tau, s, t): 0 \leq \tau \leq
s \leq t\}$, so no future values of $\Phi$ contribute.
\end{proof}

Under suitable conditions (linearization, or $F_1, F_2$ weakly nonlinear),
the composite operator admits an effective kernel approximation:
\[
  \mathcal{T}_{K_1}\circ\mathcal{T}_{K_2} \approx \mathcal{T}_{K_{\mathrm{eff}}},
  \quad K_{\mathrm{eff}} \approx K_1 \circ K_2.
\]
This provides an approximate closure of nonlinear synthesis within the kernel
algebra, connecting it back to the convolution monoid.

Spherepop specifies the categorical structure of events (morphisms,
composition, tensor product, trace). KES provides the analytic mechanism
by which those events arise from trajectory data.

\begin{definition}[KES-Spherepop functor]
$\Psi \colon \mathcal{A}_{\mathrm{KES}} \to \mathbf{Sph}$ sends each
synthesis kernel to the event morphism it generates:
$\Psi(K_{\mathrm{syn}}) = (\text{trajectories at }\mathcal{E}_\theta)$.
\end{definition}

The categories are related: KES generates the events that Spherepop
organizes.

\paragraph{Bridge to Synchronization and Quantum Structure.}
KES shows how events emerge from trajectory accumulation but does not
impose global constraints on how they combine across space. TARTAN provides
this constraint by enforcing compatibility across a tiled domain. When
local transition laws are required to be unistochastic and globally
synchronized, the resulting dynamics acquire quantum structure. Ising
synchronization selects a globally coherent phase configuration,
transforming local kernel dynamics into consistent unitary evolution.
Quantum behavior emerges as a global constraint on kernel composition.

%%=============================================================
\section{TARTAN: Recursive Tiling, Ising Synchronization, and\\Quantum Dynamics}
\label{sec:tartan}
%%=============================================================

\subsection{Overview}

TARTAN is a framework for recursive tiling and gluing of local dynamical
data. In the kernel-operator setting, TARTAN supplies the cover structure
$\{U_a\}$ over which local trajectory operators are defined, together with
the compatibility constraints that govern how they glue. The connection to
quantum mechanics arises when the local transition laws are required to be
\emph{unistochastic}: stochastic matrices arising as squared moduli of
unitary matrices. Ising synchronization then selects a globally consistent
unitary lift.

\subsection{Local Configuration Tiles}

Let $\Omega$ be covered by TARTAN tiles $\{U_a\}_{a \in A}$. Each tile
carries a finite configuration space $X_a$ and a local transition law
\[
  P^{(a)} \colon X_a \times X_a \to [0,1].
\]
We require $P^{(a)}$ to be \emph{unistochastic}: there exists a unitary
matrix $U^{(a)}$ such that
\[
  P^{(a)}_{ij} = \left|U^{(a)}_{ij}\right|^2.
\]
The matrix $P^{(a)}$ gives observable transition probabilities; $U^{(a)}$
carries the phase structure required for quantum interference.

\subsection{Phase-Lifted TARTAN Variables}

A unistochastic matrix determines probabilities but not a unique unitary
lift. We introduce local phase variables $\theta^{(a)}_{ij} \in S^1$ and write
\[
  U^{(a)}_{ij} = \sqrt{P^{(a)}_{ij}}\, e^{i\theta^{(a)}_{ij}}.
\]
The TARTAN problem is to synchronize phase lifts across overlapping tiles.
On overlaps $U_a \cap U_b$, compatibility requires:
\begin{itemize}
  \item Probability agreement: $P^{(a)}|_{U_a\cap U_b} = P^{(b)}|_{U_a\cap U_b}$.
  \item Phase coherence: $\theta^{(a)}_{ij} - \theta^{(b)}_{ij} \equiv 0
    \pmod{2\pi}$ up to admissible gauge transformations.
\end{itemize}

\subsection{Ising Synchronization}

To drive synchronization, assign an Ising spin $\sigma_a \in \{-1,+1\}$
to each tile. Neighboring tiles interact through the energy functional
\[
  E_{\mathrm{Ising}}(\sigma) = -\sum_{\langle a,b\rangle} J_{ab}\sigma_a\sigma_b
  - \sum_a h_a\sigma_a,
\]
where the coupling
\[
  J_{ab} = \exp\!\left(-\eta\left\|\theta^{(a)} - \theta^{(b)}\right\|^2_{U_a\cap U_b}\right)
\]
rewards phase-compatible tile pairs, and $h_a$ encodes local bias from
the TARTAN annotation layer. The synchronized configuration is
\[
  \sigma^* = \arg\min_\sigma E_{\mathrm{Ising}}(\sigma).
\]

\subsection{Synchronized Unistochastic Law}

Once synchronized, the local matrices glue into an effective global law:
\[
  P^{\mathrm{eff}} = \operatorname{Glue}_{\sigma^*}\!\left(\{P^{(a)}\}_{a\in A}\right),
\]
with unitary lift
\[
  U^{\mathrm{eff}}_{ij} = \sqrt{P^{\mathrm{eff}}_{ij}}\, e^{i\theta^{\mathrm{eff}}_{ij}}.
\]
The global quantum dynamics are then generated by
\[
  \psi(t+\Delta t) = U^{\mathrm{eff}}\psi(t),
  \qquad
  p_i(t+\Delta t) = \sum_j P^{\mathrm{eff}}_{ij}\, p_j(t).
\]

\subsection{Constraint Closure in the TARTAN Layer}

A globally admissible quantum process is a fixed point of the combined
TARTAN--Ising--unistochastic operator:
\[
  P^* = \mathcal{G}_{\mathrm{TARTAN}}\!\left(
    \operatorname{Sync}_{\mathrm{Ising}}\!\left(\{P^{(a)}\}\right)\right).
\]
This is an instance of the general constraint-closure pattern
$\Phi = \TK[\Phi]$, with the kernel encoding the gluing and synchronization
operations. The correspondences are:

\medskip
\begin{center}
\renewcommand{\arraystretch}{1.3}
\begin{tabular}{ll}
\toprule
\textbf{TARTAN layer} & \textbf{Kernel-operator language} \\
\midrule
Tile compatibility & Sheaf gluing condition \\
Ising alignment & Selection of consistent global section \\
TARTAN gluing & Fixed point of $\mathcal{T}_{K_{\mathrm{TARTAN}}}$ \\
Quantum transition law & Globally synchronized section $P^*$ \\
\bottomrule
\end{tabular}
\end{center}
\medskip

\begin{remark}[Core slogan]
Quantum mechanics, in this formulation, is globally synchronized unistochastic
transition dynamics: TARTAN supplies the gluing geometry; Ising synchronization
supplies the alignment dynamics; unistochasticity supplies the quantum law.
\end{remark}

\paragraph{Bridge to Symmetry Constraints.}
The TARTAN construction enforces global consistency across spatial tiles.
The next layer imposes consistency across reference frames. Where TARTAN
constrains admissibility through synchronization, relativity constrains
it through invariance of a quadratic form. Kernel operators compatible
with this symmetry must commute with Lorentz transformations, placing
geometric bounds on the operator algebra.

%%=============================================================
\section{Relativity from a Pythagorean Invariant}
\label{sec:relativity}
%%=============================================================

\subsection{The Euclidean Prototype}

In Euclidean geometry, the Pythagorean theorem defines the invariant
quadratic form $s^2 = x^2 + y^2$. Rotations preserve this quantity.
The classical square-dissection proofs demonstrate that such invariance
is equivalent to area-preserving rearrangements---the quadratic form
encodes a symmetry group.

This section argues that the Lorentz transformation and length contraction
are the exact analogues of this structure applied to spacetime, with one
sign change in the metric.

\subsection{The Spacetime Quadratic Form}

Introduce a temporal coordinate $t$ and a universal speed $c$. For a
light signal, $x = ct$, so any invariant quantity must vanish along light
trajectories. This forces the spacetime quadratic form to be
\[
  s^2 = c^2 t^2 - x^2,
\]
the Minkowski metric with indefinite signature. This is the unique
(up to scaling) quadratic form that preserves the light cone structure,
and therefore encodes the causal geometry of spacetime. It is structurally
Pythagorean, but with a sign change that encodes the distinction between
space and time.

\subsection{Invariant-Preserving Transformations}

We seek all linear maps $(x,t) \mapsto (x',t')$ preserving $s^2$:
\[
  c^2 t'^2 - x'^2 = c^2 t^2 - x^2.
\]
Assuming the transformation $x' = \gamma(x - vt)$,
$t' = \gamma(t - vx/c^2)$ and substituting into the invariant condition,
invariance holds if and only if
\[
  \gamma = \frac{1}{\sqrt{1 - v^2/c^2}}.
\]
The Lorentz factor $\gamma$ arises as the unique normalization constant
that makes the linear transformation preserve the spacetime Pythagorean form.

\begin{theorem}[Lorentz transformation from quadratic invariance]
\label{thm:lorentz}
The unique family of linear maps preserving $s^2 = c^2t^2 - x^2$ and
reducing to the identity at $v = 0$ is the Lorentz boost with parameter $\gamma$.
\end{theorem}

\begin{proof}
Substitute the general linear ansatz into $s^2 = s'^2$ and solve for the
coefficients. The constraint that $(x',t') = (x,t)$ at $v=0$ fixes the
integration constants, uniquely determining $\gamma$ as stated.
\end{proof}

\subsection{Length Contraction as Projection}

Let a rod be at rest in one frame with proper length $L_0 = x_2 - x_1$.
To measure its length in a moving frame, both endpoints must be measured at
the same time in that frame: $t'_2 = t'_1$. The condition
$t - vx/c^2 = \mathrm{const}$ combined with the Lorentz transformation for
$x'$ gives
\[
  L = x'_2 - x'_1 = \frac{L_0}{\gamma}.
\]

This is Lorentz contraction. It is not a physical compression of the rod
but a projection effect: the same invariant spacetime interval is expressed
along a different coordinate axis under a symmetry-preserving transformation.

\begin{remark}[Geometric interpretation]
In Euclidean geometry, rotation preserves circular level sets and ``dissects''
area by rearrangement. In spacetime, Lorentz boosts preserve hyperbolic level
sets of $s^2$. Length contraction is the re-slicing of those hyperbolas under
a different temporal axis---the relativistic analogue of a Euclidean dissection.
\end{remark}

\subsection{Connection to the Operator Framework}

The quadratic form $s^2$ acts as a constraint functional: it defines the
class of admissible transformations. Lorentz transformations are the symmetry
group of that constraint. In the language of kernel operators:
\[
  \text{invariant quadratic form} \;\Rightarrow\; \text{symmetry group}
  \;\Rightarrow\; \text{observable contraction effects}.
\]
This places relativity in the same structural position as the other frameworks:
a constraint on admissible dynamics, expressed through an invariant functional.

\paragraph{Bridge to Spectral Structure.}
Relativistic invariance constrains the form of admissible kernel operators
but does not explain their internal behavior: why memory induces smoothing,
why irreversibility arises, or why certain trajectories are stable. These
questions are answered by spectral analysis. The kernel operator, viewed as
an operator on path space, possesses a spectrum that governs decay,
stability, and information loss, shifting focus from external constraints
to internal dynamics.

%%=============================================================
\section{Lorentz Symmetry on Trajectory Space}
\label{sec:lorentz}
%%=============================================================

\subsection{Trajectories in Spacetime}

We extend the path space to spacetime coordinates. Let
$\Phi \colon \Omega \times [0,T] \to \R$ be a trajectory, and equip the
domain with the Minkowski quadratic form $s^2 = c^2t^2 - x^2$. A Lorentz
transformation $\Lambda$ acts on spacetime coordinates by
\[
  (x,t) \mapsto (x',t') = \Lambda(x,t), \quad s'^2 = s^2.
\]
This induces a pullback operator on trajectories:
\[
  (\Lambda\cdot\Phi)(x,t) := \Phi(\Lambda^{-1}(x,t)).
\]

\subsection{Lorentz Compatibility of Kernel Operators}

\begin{definition}[Lorentz-compatible operator]
A trajectory operator $\TK$ is \emph{Lorentz-compatible} if
\[
  \Lambda\cdot\TK[\Phi] = \TK[\Lambda\cdot\Phi]
\]
for all trajectories $\Phi$ and all Lorentz transformations $\Lambda$.
\end{definition}

\begin{proposition}[Invariant kernels are Lorentz-compatible]
\label{prop:lorentz}
If the kernel $K$ depends only on the spacetime interval, i.e.\
$K(t,\tau) = \tilde{K}(s^2)$ for some function $\tilde{K}$, then
$\TK$ is Lorentz-compatible.
\end{proposition}

\begin{proof}
Since $s^2$ is invariant under $\Lambda$, we have $K(t',\tau') = K(t,\tau)$.
Substituting into the definition of $\TK$:
\[
  (\Lambda\cdot\TK[\Phi])(x,t)
  = \TK[\Phi](\Lambda^{-1}(x,t))
  = \int_0^{t^*} K(t^*,\tau)\,F(\Phi(\Lambda^{-1}(x,\tau)))\,d\tau,
\]
where $t^* = [\Lambda^{-1}(x,t)]_t$. Under the change of variables
$\tau \mapsto \Lambda^{-1}(\tau)$, this equals $\TK[\Lambda\cdot\Phi](x,t)$,
since $K$ is invariant.
\end{proof}

\subsection{Relativistic Kernel Evolution}

A fully Lorentz-invariant trajectory operator integrates over the past
light cone rather than a fixed time interval:
\[
  \TK[\Phi](x,t)
  = \int_{\mathcal{C}(x,t)} K\bigl(s^2((x,t),(x',t'))\bigr)\,
    F(\Phi(x',t'))\, d\mu(x',t'),
\]
where $\mathcal{C}(x,t) = \{(x',t') : c^2(t-t')^2 - (x-x')^2 \geq 0,\;
t' \leq t\}$ is the past light cone. This is the relativistic generalization
of the causal Volterra operator: instead of integrating over $[0,t]$, it
integrates over all causally accessible past events.

\subsection{Length Contraction as Trajectory Projection}

Let $\Phi$ represent a spatial configuration at rest in one frame.
Different inertial frames correspond to different temporal slicings of
spacetime. Since $\Lambda$ preserves $s^2$, the underlying trajectory is
unchanged; what changes is its projection onto the spatial axis of the
moving frame. This projection is exactly Lorentz contraction
$L = L_0/\gamma$: a consequence of evaluating the same invariant trajectory
under different coordinate projections.

\subsection{Interpretation}

Within the kernel-operator framework, Lorentz symmetry is a compatibility
condition rather than an independent postulate. The diagram
\[
  \Lambda\cdot\TK = \TK\cdot\Lambda
\]
says that temporal evolution and frame change commute---a statement about
the algebraic structure of $\Kcal\mathsf{Op}$, not a separate postulate
of physics. Relativity appears as the constraint that kernel operators
must commute with the symmetry group of the invariant quadratic form.

%%=============================================================
\section{Constraint Closure and the Trajectory-First Theory}
\label{sec:constraint}
%%=============================================================

\subsection{Trajectories as Primary Objects}

The standard picture: state space $\mathcal{X}$, evolution map $\Phi_t
\colon \mathcal{X} \to \mathcal{X}$. Completion = reaching a fixed point
$x^* \in \mathcal{X}$.

The trajectory-first picture: the primary objects are trajectories
$f \in \Pcal$, and states are derived by evaluation. Completion is not
``this state is fixed'' but:

\begin{center}
\textit{this trajectory satisfies $\mathcal{T}_K[f^*] = f^*$.}
\end{center}

This is the fixed-point condition for the trajectory operator $\TK$.

\subsection{Constraint Operators}

\begin{definition}[Constraint operator]
A \emph{constraint operator} on $\Pcal$ is $\mathcal{T}_K$ satisfying:
\begin{enumerate}[label=(\roman*)]
  \item \textbf{Causal}: $\mathcal{T}_K[f](t)$ depends only on $f|_{[0,t]}$.
  \item \textbf{Contractive}: $\|\mathcal{T}_K[f] - \mathcal{T}_K[g]\|_\Pcal
    \leq L\|f-g\|_\Pcal$ for some $L < 1$.
  \item \textbf{Positivity-preserving}: $f \geq 0 \Rightarrow
    \mathcal{T}_K[f] \geq 0$.
\end{enumerate}
\end{definition}

\begin{definition}[Constraint closure]
A trajectory $f^*$ is \emph{closed under constraint $\mathcal{T}_K$} if
$\mathcal{T}_K[f^*] = f^*$: it is a fixed point of its own constraint.
\end{definition}

\subsection{The Constraint Closure Theorem}

\begin{theorem}[Constraint Closure Theorem]
\label{thm:closure}
Let $\mathcal{T}_K$ be a constraint operator. Then:
\begin{enumerate}[label=(\roman*)]
  \item There exists a unique fixed point $f^* \in \Pcal$ with
    $\mathcal{T}_K[f^*] = f^*$.
  \item For any $f_0 \in \Pcal$, the iterates $f_{n+1} = \mathcal{T}_K[f_n]$
    converge to $f^*$ in $\Pcal$-norm.
  \item The operator $\mathcal{T}_K$ is idempotent on $\{f^*\}$:
    $\mathcal{T}_K(\mathcal{T}_K[f^*]) = \mathcal{T}_K[f^*]$.
\end{enumerate}
\end{theorem}

\begin{proof}
Parts (i) and (ii): the Banach fixed-point theorem applied to the contractive
map $\mathcal{T}_K$ on the complete metric space $(\Pcal, \|\cdot\|)$.
Part (iii): $f^* = \mathcal{T}_K[f^*]$ implies $\mathcal{T}_K(f^*) =
f^* = \mathcal{T}_K[f^*]$, so applying $\mathcal{T}_K$ again gives the
same result.
\end{proof}

\subsection{Idempotence as Completion}

\begin{corollary}[Completion without states]
A process defined by trajectory $f \in \Pcal$ and constraint $\mathcal{T}_K$
is \emph{complete} if and only if $\mathcal{T}_K$ is idempotent on $f$.
No state space is required.
\end{corollary}

This is the formal statement of the trajectory-first completion principle:
\[
  \text{completion} = \text{constraint closure} = \text{kernel idempotence}.
\]

\subsection{Yarncrawler as Constraint Closure}

The Yarncrawler kernel $K_Y$ is idempotent by Proposition~\ref{prop:idempotent}.
This is not a coincidence: Yarncrawler reconstruction is exactly constraint
closure for the sheaf-consistency constraint.

\begin{proposition}[Yarncrawler is a constraint closure]
Yarncrawler reconstruction is an instance of Theorem~\ref{thm:closure}
with $\mathcal{T}_K = \mathcal{T}_{K_Y}$.
\end{proposition}

The world-state reconstruction completes when the Yarncrawler constraint
closes: when further application of $\mathcal{T}_{K_Y}$ produces no new
information.

\paragraph{Bridge to Spectral and Variational Analysis.}
Constraint closure establishes when a trajectory is a fixed point. Spectral
theory now reveals why such fixed points are stable and what happens to
trajectories that do not satisfy the constraint: their non-invariant
components are spectrally damped. Variational theory shows that the same
fixed point is a minimizer of a mismatch functional. The following two
sections formalize this triple equivalence before the Unified Memory Theorem
assembles it into a single statement.


%%=============================================================
\section{Spectral Theory of Kernel Operators}
\label{sec:spectral}
%%=============================================================

\subsection{Operator-Theoretic Setting}

We now study $\TK$ as a linear operator on a function space, characterizing
memory, irreversibility, and convergence through spectral properties.
Let $\Pcal := L^2([0,T], X)$ and define the linear kernel operator
\[
  (\mathcal{K}f)(t) := \int_0^t K(t,\tau)\,f(\tau)\,d\tau,
  \quad \mathcal{K} \colon \Pcal \to \Pcal.
\]
This is bounded by Lemma~\ref{lem:bounded}.

\subsection{Compactness and Smoothing}

\begin{theorem}[Compactness of causal kernel operators]
\label{thm:compact}
If $K \in L^2(\Delta_T)$, then $\mathcal{K}$ is a compact operator on $\Pcal$.
\end{theorem}

\begin{proof}
$\mathcal{K}$ is a Hilbert--Schmidt operator:
$\|\mathcal{K}\|_{HS}^2 = \int_0^T\!\int_0^t|K(t,\tau)|^2\,d\tau\,dt < \infty$.
All Hilbert--Schmidt operators are compact.
\end{proof}

\begin{remark}[Smoothing interpretation]
Compactness means $\mathcal{K}$ maps bounded sets into relatively compact
sets. Physically, this is \emph{smoothing}: fine-scale variations in the
trajectory are suppressed by integration against $K$. Memory is not merely
accumulation---it is selective persistence with suppression of high-frequency
modes.
\end{remark}

\subsection{Spectral Structure and Memory Decay}

\begin{theorem}[Spectral structure]
\label{thm:spectrum}
If $\mathcal{K}$ is compact, then $\sigma(\mathcal{K}) = \{0\}\cup\{\lambda_n\}$
with $\lambda_n \to 0$, each nonzero $\lambda_n$ has finite multiplicity, and
$0$ is the only accumulation point.
\end{theorem}

The eigenvalues $\lambda_n$ measure how strongly different modes of the
trajectory persist under the kernel. Since $\lambda_n \to 0$, all modes
eventually decay: \emph{memory is lossy}. Repeated application of the
kernel erases information.

\subsection{Irreversibility from Spectral Decay}

\begin{theorem}[Non-invertibility of smoothing kernels]
\label{thm:noninvertspectral}
If $K \in L^1(\Delta_T)$ and is not a distribution concentrated on $t = \tau$,
then $\mathcal{K}$ is not invertible on $\Pcal$.
\end{theorem}

\begin{proof}
A compact operator on an infinite-dimensional Banach space cannot be
invertible, since its spectrum accumulates at $0$.
\end{proof}

\begin{remark}[Algebraic irreversibility]
This gives a purely spectral derivation of irreversibility:
$\text{smoothing} \Rightarrow \lambda_n \to 0 \Rightarrow \text{no inverse}$.
No thermodynamic argument is required; information loss is encoded in the
operator spectrum.
\end{remark}

\subsection{Fractional Kernels and Power-Law Spectra}

\begin{proposition}[Spectral decay of fractional operators]
\label{prop:fractionalspec}
The singular values of $\mathcal{K}_\alpha$ satisfy
$s_n(\mathcal{K}_\alpha) \sim n^{-(1-\alpha)}$.
\end{proposition}

For $\alpha \to 0$: rapid spectral decay, short memory. For $\alpha \to 1$:
slow decay, long memory. The fractional parameter $\alpha$ directly controls
the spectral decay rate, making memory depth a measurable spectral quantity.

\subsection{Stability of Fixed Points}

Linearizing $\TK[f] = \mathcal{K}(F\circ f)$ around a fixed point $f^*$:

\begin{theorem}[Stability criterion]
\label{thm:stability}
A fixed point $f^*$ of $\TK$ is stable if
$\rho\bigl(\mathcal{K}\circ DF(f^*)\bigr) < 1$,
where $\rho$ is the spectral radius.
\end{theorem}

Constraint closure is therefore not just existence of a fixed point but
\emph{spectral contraction} of deviations: the kernel controls stability
by damping perturbations.

\subsection{Resolvent and Memory Propagation}

The resolvent $R(\lambda) = (I - \lambda\mathcal{K})^{-1}$ expands as
\[
  R(\lambda) = \sum_{n=0}^\infty \lambda^n \mathcal{K}^n,
  \quad |\lambda| < 1/\|\mathcal{K}\|.
\]
Each term represents a deeper layer of historical influence. Memory
propagation is literally a geometric series in kernel composition.

\paragraph{Bridge to Variational Equivalence.}
Spectral theory characterizes the behavior of $\mathcal{K}$ through
eigenstructure. The same structure can be obtained from an optimization
principle: fixed points of $\mathcal{K}$ are minimizers of a natural
mismatch functional. The next section formalizes this equivalence,
completing the identification between operator algebra, energy minimization,
and constraint closure.

%%=============================================================
\section{Spectral--Variational Equivalence}
\label{sec:spectral_variational}
%%=============================================================

\subsection{Overview}

We establish the equivalence between three perspectives on $\TK$:

\begin{center}
\textit{variational minimization}
$\;\Longleftrightarrow\;$
\textit{spectral contraction}
$\;\Longleftrightarrow\;$
\textit{constraint closure}.
\end{center}

\subsection{The Mismatch Functional}

\begin{definition}[Mismatch functional]
\[
  \mathcal{J}[f] := \tfrac{1}{2}\|f - \mathcal{K}f\|^2_{\Pcal}.
\]
\end{definition}

\begin{proposition}[Euler--Lagrange condition]
\label{prop:EL}
A trajectory $f^*$ is a stationary point of $\mathcal{J}$ if and only if
$(I - \mathcal{K})^*(I - \mathcal{K})\,f^* = 0$.
\end{proposition}

\begin{proof}
$\delta\mathcal{J} = \langle (I-\mathcal{K})f,\,(I-\mathcal{K})\delta f\rangle = 0$
for all $\delta f$ gives $(I-\mathcal{K})^*(I-\mathcal{K})f = 0$.
\end{proof}

\begin{corollary}[Variational constraint closure]
$\mathcal{J}[f^*] = 0 \Longleftrightarrow f^* = \mathcal{K}f^*$.
\end{corollary}

Minimizing $\mathcal{J}$ enforces the fixed-point condition: constraint
closure is a variational principle.

\subsection{Spectral Characterization of Minimizers}

Expand $f = \sum_n c_n e_n$ in the eigenbasis of $\mathcal{K}$.
Then $\mathcal{J}[f] = \tfrac{1}{2}\sum_n(1-\lambda_n)^2|c_n|^2$.

\begin{proposition}[Spectral minimizers]
$f^*$ minimizes $\mathcal{J}$ if and only if $c_n = 0$ whenever $\lambda_n \neq 1$.
The fixed-point space is the eigenspace $\ker(I - \mathcal{K})$.
\end{proposition}

\subsection{Gradient Flow Convergence}

The gradient flow $\dot{f} = -\nabla\mathcal{J}[f] = -(I-\mathcal{K})^*(I-\mathcal{K})f$
diagonalizes as $\dot{c}_n = -(1-\lambda_n)^2 c_n$.

\begin{theorem}[Convergence to fixed point]
\label{thm:gradflow}
If $\rho(\mathcal{K}) < 1$ on the complement of $\ker(I-\mathcal{K})$, then
$f(t) \to f^* \in \ker(I-\mathcal{K})$.
\end{theorem}

\begin{proof}
For $\lambda_n \neq 1$: $c_n(t) \to 0$ exponentially. For $\lambda_n = 1$:
$c_n$ is constant. The system projects onto the constraint-closed subspace.
\end{proof}

\subsection{The Equivalence Theorem}

\begin{theorem}[Spectral--Variational Equivalence]
\label{thm:sve}
For a compact kernel operator $\mathcal{K}$, the following are equivalent:
\begin{enumerate}[label=(\roman*)]
\item $f^*$ minimizes $\mathcal{J}[f] = \tfrac{1}{2}\|f - \mathcal{K}f\|^2$,
\item $f^*$ satisfies $\mathcal{K}f^* = f^*$,
\item $f^*$ lies in the eigenspace $\lambda = 1$ of $\mathcal{K}$,
\item $f^*$ is the limit of the gradient flow of $\mathcal{J}$,
\item $f^*$ is stable under iteration of $\mathcal{K}$.
\end{enumerate}
\end{theorem}

\begin{proof}
(i)$\Rightarrow$(ii): Corollary above.
(ii)$\Rightarrow$(iii): definition of eigenvector.
(iii)$\Rightarrow$(iv): Theorem~\ref{thm:gradflow}.
(iv)$\Rightarrow$(v): fixed point implies invariance under iteration.
(v)$\Rightarrow$(ii): the iteration limit satisfies $\mathcal{K}f^* = f^*$.
\end{proof}

\begin{center}
\textit{Constraint closure is projection onto the unit-eigenspace of a
smoothing operator.}
\end{center}

\paragraph{Bridge to Worked Example.}
The preceding sections establish that kernel dynamics, spectral decay,
variational minimization, and constraint closure are equivalent at the
level of abstract structure. To make this equivalence concrete and
eliminate any remaining separation between perspectives, we now exhibit
a single system---fractional diffusion---in which all of them coincide
explicitly.

%%=============================================================
\section{Unified Worked Example: Fractional Diffusion}
\label{sec:worked}
%%=============================================================

\subsection{Overview}

We carry fractional diffusion through every representation developed in
this monograph. The goal is structural: to show that kernel dynamics,
spectral decay, variational minimization, constraint closure, RSVP field
evolution, and sheaf consistency are the same object described in different
languages.

\subsection{Classical vs.\ Fractional Formulation}

Classical diffusion: $\partial_t u(x,t) = D\Delta u(x,t)$.
Local and Markovian.

Fractional diffusion (Caputo): $D_t^\alpha u = D\Delta u$, $\alpha \in (0,1)$.
Expanding:
\[
  \frac{1}{\Gamma(1-\alpha)}\int_0^t (t-\tau)^{-\alpha}\partial_\tau u(x,\tau)\,d\tau
  = D\Delta u(x,t).
\]
This is $\mathcal{T}_{K_\alpha}[\partial_t u](t) = D\Delta u(t)$, with
$K_\alpha(t,\tau) = (t-\tau)^{-\alpha}/\Gamma(1-\alpha)$.

\subsection{Kernel Operator Form and Fixed Point}

Rewriting as a Volterra equation:
\[
  u(x,t) = u_0(x) + \int_0^t K_\alpha(t,\tau)\,D\Delta u(x,\tau)\,d\tau.
\]
This is $u = u_0 + \mathcal{K}[u]$, so solutions are fixed points of the
trajectory operator $u \mapsto u_0 + \mathcal{T}_{K_\alpha}[u]$.

\subsection{Spectral Representation}

Expand over the Laplacian eigenbasis $\Delta\phi_n = -\mu_n\phi_n$.
Each mode satisfies $D_t^\alpha c_n = -D\mu_n c_n$, with solution
\[
  c_n(t) = c_n(0)\,E_\alpha(-D\mu_n t^\alpha),
\]
where $E_\alpha$ is the Mittag--Leffler function. For large $t$:
$c_n(t) \sim t^{-\alpha}$ (power-law decay), contrasting with classical
exponential decay. \emph{Memory appears as slow spectral decay.}

\subsection{Variational Formulation}

Define
$\mathcal{J}_\alpha[u] = \tfrac{1}{2}\int_0^T\!\|\mathcal{T}_{K_\alpha}[\partial_t u] - D\Delta u\|^2\,dt$.
Minimizers of $\mathcal{J}_\alpha$ satisfy the fractional diffusion equation.
Thus: \emph{diffusion = variational minimization of kernel mismatch}.

\subsection{RSVP Interpretation}

Set $\Phi = u$, $\mathbf{v} = 0$, $S = 0$ in the kernel-weighted RSVP
system. Equations~\eqref{eq:kphi}--\eqref{eq:kS} reduce to
$\partial_t\Phi = \mathcal{T}_{K_\alpha}[D\Delta\Phi]$.
Fractional diffusion is the \emph{minimal RSVP system}: scalar field only,
no transport, no entropy coupling, pure memory-weighted diffusion.

\subsection{KES Interpretation}

Define synthesis operator $\mathcal{T}_{K_{\mathrm{syn}}}[u](t) = \int_0^t
\psi(t)K_\alpha(t,\tau)u(\tau)\,d\tau$. Events occur when
$|\mathcal{T}_{K_{\mathrm{syn}}}[u](t)| \geq \theta$. Continuous diffusion
generates discrete events as threshold crossings of smoothed trajectories.

\subsection{Yarncrawler Interpretation}

Local solutions $\{u_i\}$ on a cover $\{U_i\}$ satisfy compatibility on
overlaps. The global solution exists iff $\check{H}^1 = 0$. Fractional
diffusion solutions are global sections of the trajectory sheaf.

\subsection{Spectral--Variational--Closure Summary}

\medskip
\begin{center}
\renewcommand{\arraystretch}{1.4}
\begin{tabular}{ll}
\toprule
\textbf{Perspective} & \textbf{Statement for fractional diffusion} \\
\midrule
Kernel & $u = u_0 + \mathcal{T}_{K_\alpha}[u]$ \\
Spectral & $c_n(t) \sim E_\alpha(-\mu_n t^\alpha) \sim t^{-\alpha}$ \\
Variational & $u = \arg\min\,\mathcal{J}_\alpha$ \\
Constraint closure & $u$ is a fixed point of $u \mapsto u_0 + \mathcal{K}[u]$ \\
RSVP & Minimal kernel-weighted scalar field \\
Sheaf & Global section of trajectory sheaf \\
\bottomrule
\end{tabular}
\end{center}
\medskip

\begin{center}
\textit{All descriptions produce the same solution space.
Diffusion is not a differential equation; it is a fixed point
of a kernel operator whose stability is spectral, whose structure
is variational, and whose behavior is memory.}
\end{center}

%%=============================================================
\section{The Unified Memory Theorem}
\label{sec:unified}
%%=============================================================

\subsection{Statement}

We now assemble all eight frameworks into a single categorical statement.

\begin{theorem}[Unified Memory Theorem]
\label{thm:unified}
There is a commutative diagram of functors
\[
\begin{array}{ccccc}
  \mathcal{A}_{\mathrm{KES}}
    & \xrightarrow{\;\Psi\;}
  & \mathbf{Sph}
    & \xrightarrow{\;\Omega\;}
  & \Kcal\mathsf{Op} \\
  & & \downarrow\scriptstyle{\Lambda}
    & & \downarrow\scriptstyle{\mathcal{R}} \\
  & & \mathbf{Yarn}
    & \xrightarrow{\;\Xi\;}
  & \mathbf{FieldSys}
\end{array}
\]
where $\Psi$ is the KES-Spherepop functor, $\Omega$ embeds events via
kernel representation, $\mathcal{R}$ is the RSVP functor, $\Lambda$ sends
events to their post-event trajectories, $\mathbf{Yarn}$ is the Yarncrawler
category of sheaf sections, and $\Xi$ embeds global sections as RSVP field
histories.

Moreover, the following conditions on a trajectory $\Phi \in \Pcal$ are
equivalent characterizations of admissible trajectory realization:
\begin{enumerate}[label=(\roman*)]
  \item \textbf{Constraint closure}: $\Phi = \mathcal{T}_{\mathrm{tot}}[\Phi]$
    (fixed point of the synthesized kernel operator).
  \item \textbf{Sheaf consistency}: $\Phi$ is a global section of the sheaf
    of local kernel solutions $\{\Phi_i\}$ over a cover $\{U_i\}$.
  \item \textbf{Event fixed point}: $\Phi$ is a fixed point of a traced
    composition of irreversible Spherepop events.
  \item \textbf{Kernel algebra}: $\Phi$ is a fixed point of a composition
    in $\Kcal\mathsf{Op}$.
  \item \textbf{Symmetry compatibility}: $\Phi$ is stable under the symmetry
    group of the spacetime invariant $s^2 = c^2t^2 - x^2$, i.e.\
    $\Lambda\cdot\mathcal{T}_{\mathrm{tot}}[\Phi] =
    \mathcal{T}_{\mathrm{tot}}[\Lambda\cdot\Phi]$.
\end{enumerate}

In addition:
\begin{enumerate}[label=(\alph*)]
  \item Every RSVP field solution is the image under $\Xi$ of a Yarncrawler
    global section.
  \item Every Yarncrawler reconstruction corresponds to an idempotent
    constraint operator in $\Kcal\mathsf{Op}$.
  \item Every Spherepop event is a kernel morphism; its trace is
    irreversible (Theorem~\ref{thm:irrev}).
  \item Every KES event is a threshold-crossing of
    $\mathcal{T}_{K_{\mathrm{syn}}}$ (Proposition~\ref{prop:discrete}).
  \item TARTAN quantum dynamics arise as globally synchronized unistochastic
    sections, which are fixed points of the gluing operator
    $\mathcal{G}_{\mathrm{TARTAN}}\circ\operatorname{Sync}_{\mathrm{Ising}}$.
  \item Lorentz-compatible kernel operators commute with the symmetry group
    of $s^2$ (Proposition~\ref{prop:lorentz}).
\end{enumerate}
\end{theorem}

\begin{proof}
\textbf{Commutativity of the diagram.} The two paths from $\mathbf{Sph}$
to $\mathbf{FieldSys}$ are: (1) embed as kernel, apply RSVP functor;
(2) extract post-event trajectory, embed as global sheaf section, embed
in field history. Both paths yield the same RSVP history because both
solve the same kernel-weighted Volterra equation (Definition~\ref{def:krsvp})
with initial data determined by the event $e$.

\textbf{(i)$\Leftrightarrow$(ii).} A fixed point of $\mathcal{T}_{\mathrm{tot}}$
defines a global solution whose restrictions to each $U_i$ are local solutions.
Conversely, compatible local solutions glue (Theorem~\ref{thm:sheafvar}) to a
global section that satisfies $\Phi = \mathcal{T}_{\mathrm{tot}}[\Phi]$.

\textbf{(i)$\Leftrightarrow$(iii).} Event compositions correspond to kernel
compositions (Proposition~\ref{prop:totalkernel}); fixed points of the
resulting operator are trajectories invariant under the composed event.
Traced events yield the fixed-point equation $\Phi = E(\Phi)$, which is
precisely constraint closure.

\textbf{(i)$\Leftrightarrow$(iv).} Kernel composition in $\Kcal\mathsf{Op}$
defines the same algebra as trajectory synthesis; fixed points are the
same in both descriptions.

\textbf{(i)$\Leftrightarrow$(v).} For Lorentz-compatible kernels
(Proposition~\ref{prop:lorentz}), fixed points of $\mathcal{T}_{\mathrm{tot}}$
are automatically Lorentz-stable. Conversely, Lorentz-stable trajectories
that are solutions of the kernel equation are fixed points.

\textbf{(a)--(d)} were established in the individual framework sections
above.

\textbf{(e)} TARTAN gluing defines a sheaf over the tile cover; Ising
synchronization selects a globally consistent phase section; the result is
a fixed point of the composed gluing operator, which is an instance of (i).

\textbf{(f)} Proposition~\ref{prop:lorentz}.
\end{proof}

\subsection{The Eight-Point Translation}

The Unified Memory Theorem says in words: memory, causality, irreversibility,
global consistency, constraint closure, quantum synchronization, and
relativistic invariance are perspectives on a single structure. The
translation table:

\medskip
\begin{center}
\renewcommand{\arraystretch}{1.4}
\begin{tabular}{ll}
\toprule
\textbf{Property / Framework} & \textbf{Algebraic form in $\Kcal\mathsf{Op}$} \\
\midrule
Memory & Non-locality of $K(t,\tau)$ in time \\
Causality & Triangular support: $K(t,\tau) = 0$ for $\tau > t$ \\
Irreversibility & Non-invertibility of the trace morphism \\
Global consistency & Vanishing of $\check{H}^1(\{U_i\}, \mathscr{T})$ \\
Constraint closure & Idempotence / fixed point of $\mathcal{T}_K$ \\
TARTAN quantum law & Fixed point of $\mathcal{G}_{\mathrm{TARTAN}}\circ\operatorname{Sync}_{\mathrm{Ising}}$ \\
Lorentz symmetry & Commutativity $\Lambda\cdot\TK = \TK\cdot\Lambda$ \\
Fractional interpolation & Semigroup: $K_\alpha\circ K_\beta = K_{\alpha+\beta}$ \\
\bottomrule
\end{tabular}
\end{center}
\medskip

Each entry is a theorem about kernel operators, not an assumption about
the physical world.

\subsection{The Self-Modulating Kernel: A Future Direction}

The monograph has held the kernel $K$ fixed. A natural generalization---noted
but deliberately deferred in the notes---is to let $K$ depend on the
system's own state:
\[
  K = K(t,\tau;\Phi(t),\mathbf{v}(t),S(t)).
\]
In this regime:
\begin{itemize}
  \item memory depth depends on entropy: high-entropy regions forget faster;
  \item transport reshapes its own history weighting;
  \item the system controls how strongly it remembers.
\end{itemize}

This is no longer fractional calculus in the classical sense. It is a
self-modulating nonlocal field theory in which $K$ is itself a dynamical
variable. The Unified Memory Theorem applies to the fixed-kernel case; the
self-modulating case is the natural continuation of this program.

%%=============================================================
\section{Discussion}
\label{sec:discussion}
%%=============================================================

\subsection{Memory as the Default}

The standard approach treats memory as a complication: an extra ingredient
added to explain anomalous behavior. The kernel-operator framework inverts
this. The delta kernel ($K = \delta$) is the \emph{exceptional} case:
it is the only kernel with zero support width. Every other kernel has
positive support and therefore encodes memory. The generic situation is
history-dependent dynamics; Markovian dynamics are the measure-zero limit.

The twentieth-century emphasis on Markovian models was an architectural
choice, not an empirical finding. Locality was imposed because it makes
differential equations tractable. The kernel-operator framework recovers
the generality that locality sacrifices, without abandoning rigor.

\subsection{The Formula Sheet, Revisited}

The opening observation of the monograph was that a standard formula sheet
is an atlas of local charts. The conclusion follows: \textit{the formula
sheet is a discrete sampling of a continuous operator space whose coordinates
are kernel choice $K$, operator order $\alpha$, and constraint structure $F$.}

Arithmetic identities are charts near $K = \delta$, $\alpha = 0$.
Integration formulas are charts near $K = 1$, $\alpha = 1$.
Fractional operators are the curves connecting them.
RSVP field equations are surfaces in the space parameterized by
$(\alpha, \beta, \gamma)$.
Yarncrawler, Spherepop, KES, and constraint closure are structured
submanifolds of the same space.

The unified picture does not render the formula sheet obsolete; it
reveals the hidden topology of the space from which each formula was
sampled.

\subsection{Irreversibility from Algebra, Not from Statistics}

A long-standing puzzle: how does irreversibility emerge from time-reversible
microscopic laws? The standard resolution invokes thermodynamic limits,
coarse-graining, or cosmological initial conditions.

Theorem~\ref{thm:irrev} offers an algebraic perspective: irreversibility
is a property of the \emph{trace operation} in the symmetric monoidal
category of events, grounded in the entropy growth of the RSVP embedding.
No thermodynamic limit is required. The non-invertibility of the traced
morphism is a consequence of operator structure.

This does not resolve the foundational puzzle of statistical mechanics
(which concerns Hamiltonian dynamics specifically), but it suggests
that irreversibility can be located in the morphism structure of a physical
theory rather than in statistical approximations.

\subsection{Quantum Mechanics as Synchronized Trajectory Dynamics}

The TARTAN section establishes a structural connection between kernel-operator
theory and quantum mechanics via Barandes's unistochastic reformulation. The
key insight is that quantum transition probabilities arise as squared moduli
of unitary matrices---a constraint (unistochasticity) that, when enforced
locally and synchronized globally via Ising dynamics, produces a global quantum
law. This is an instance of the general TARTAN--sheaf pattern: local sections
(tile transition laws) satisfying compatibility conditions (phase coherence)
glue to a global section (the effective quantum dynamics).

The implication is that quantum behavior need not be posited as a primitive.
It emerges from the same local-to-global constraint structure that governs
RSVP field evolution and Yarncrawler reconstruction. The trajectory operator
$\TK$ with the TARTAN kernel encodes both the probabilistic structure (through
unistochasticity) and the memory structure (through the fractional kernel
weighting).

\subsection{Relativity as Symmetry Constraint on Kernels}

Section~\ref{sec:relativity} showed that Lorentz transformations arise from
a single requirement: that a quadratic form $s^2 = c^2t^2 - x^2$ be invariant.
Section~\ref{sec:lorentz} then embedded this into the kernel-operator framework:
Lorentz-compatible kernels are those that depend only on invariant spacetime
intervals, and the commutativity condition $\Lambda\cdot\TK = \TK\cdot\Lambda$
is the algebraic statement of relativistic invariance.

This means that the transition from non-relativistic to relativistic field
theory is, in the kernel-operator language, a transition from
arbitrary causal kernels to kernels that are invariant under the Lorentz group.
It is a restriction on the kernel, not an independent postulate about the world.


Constraint closure has a cognitive reading: a thinking system operating on
trajectories (sequences of states, representations, or actions) completes
a task when its constraint operator becomes idempotent. This is the algebraic
definition of ``done.''

The trajectory-first perspective contrasts with goal-state architectures,
in which completion is defined by reaching a designated terminal state.
In the constraint-closure framework, completion \emph{emerges} from the
dynamics of the system---the fixed point is not specified in advance but
reached by iteration. This has implications for autonomous systems designed
to complete open-ended tasks without explicit terminal-state specifications.

The connection to representation learning: kernel-weighted prediction
(which uses the full trajectory history, weighted by $K_\alpha$) generalizes
the Markovian prediction of current transformer architectures. The semigroup
law provides a principled interpolation between these poles parameterized
by $\alpha$.

\subsection{Open Problems}

\begin{enumerate}
  \item \textbf{Self-modulating kernels.} Formalize the dynamics of
    $K(t,\tau;\Phi,\mathbf{v},S)$ as a field equation in its own right.
    What is the stability theory?

  \item \textbf{Kernel renormalization.} Do the orders $(\alpha,\beta,\gamma)$
    of the RSVP system flow under renormalization-group transformations?
    What are the fixed points?

  \item \textbf{Computability of constraint closure.} Given a computable
    $\mathcal{T}_K$, is the fixed point $f^*$ always computable?

  \item \textbf{Sheaf cohomology and information geometry.} Is there a
    natural Fisher metric on the space of sheaf sections making $\mathcal{E}$
    an information-geometric functional?

  \item \textbf{KES phase transitions.} What happens at the exceptional
    thresholds excluded by Proposition~\ref{prop:discrete}? Do they
    correspond to phase transitions in the event structure?

  \item \textbf{Quantum extension.} Do the semigroup law and Unified Memory
    Theorem extend to operator-valued kernels on Hilbert spaces?
\end{enumerate}

%%=============================================================
\section{Conclusion}
\label{sec:conclusion}
%%=============================================================

This monograph has developed a single thesis: that dynamics are not
functions of state but functionals of trajectories, and that kernel
operators provide the minimal algebra for expressing that dependence.

The argument proceeded through five moves and eight frameworks:

\begin{enumerate}
  \item \textbf{Reframe the formula sheet.} Standard mathematical formulas
    are local charts in a continuous operator space. The hidden assumption
    underlying all differential and integral formulas is locality: the
    kernel $K = \delta$.

  \item \textbf{Introduce the trajectory operator.} The central object
    $\TK[f](t) = \int_0^t K(t,\tau) F(f(\tau))\, d\tau$ unifies
    integration, differentiation, fractional calculus, RSVP transport,
    KES synthesis, and constraint closure as specializations of a single
    form.

  \item \textbf{Prove the semigroup law.} The fractional kernel family
    $\{K_\alpha\}$ is a one-parameter semigroup under composition. This
    is a consequence of the Beta function identity, not a definition.
    It makes fractional calculus algebraically inevitable.

  \item \textbf{Categorify.} The trajectory operators form $\Kcal\mathsf{Op}$
    with $\mathbf{Ker}$ as its linear subcategory. The original five
    frameworks---RSVP, Yarncrawler, Spherepop, KES, Constraint Closure---are
    subcategories and functorial images. TARTAN quantum dynamics emerge as
    globally synchronized sections; Lorentz symmetry appears as a
    commutativity condition on kernel operators.

  \item \textbf{Prove the Unified Memory Theorem.} The commutative
    diagram of Theorem~\ref{thm:unified} establishes that all eight
    frameworks are not independent theories but equivalent characterizations
    of a single condition: a trajectory is a fixed point of its synthesized
    kernel operator.
\end{enumerate}

The result is a picture in which memory is not a property added to a
dynamics but the defining characteristic of dynamics as such. Systems
that remember their history are the rule; systems that forget are the
exception produced by setting $K = \delta$.

The formula sheet, revisited: it is not a list but a sample. Every
formula in it is a point in the continuous operator space parameterized
by kernel, order, and constraint structure. Fractional calculus is the
curve connecting two such points. RSVP, Yarncrawler, Spherepop, KES,
Constraint Closure, TARTAN, and relativistic symmetry are structured
regions of the same space.

The algebra of kernels is the algebra of memory. Memory is the algebra
of time.

\begin{enumerate}
  \item \textbf{Reframe the formula sheet.} Standard mathematical formulas
    are local charts in a continuous operator space. The hidden assumption
    underlying all differential and integral formulas is locality: the
    kernel $K = \delta$.

  \item \textbf{Introduce the trajectory operator.} The central object
    $\TK[f](t) = \int_0^t K(t,\tau) F(f(\tau))\, d\tau$ unifies
    integration, differentiation, fractional calculus, RSVP transport,
    KES synthesis, and constraint closure as specializations of a single
    form.

  \item \textbf{Prove the semigroup law.} The fractional kernel family
    $\{K_\alpha\}$ is a one-parameter semigroup under composition. This
    is a consequence of the Beta function identity, not a definition.
    It makes fractional calculus algebraically inevitable.

  \item \textbf{Categorify.} The trajectory operators form the category
    $\Kcal\mathsf{Op}$, in which the semigroup becomes a functor. The
    five frameworks---RSVP, Yarncrawler, Spherepop, KES, Constraint
    Closure---are subcategories and functorial images.

  \item \textbf{Prove the Unified Memory Theorem.} The commutative
    diagram of Theorem~\ref{thm:unified} establishes that the five
    frameworks are not independent theories but linked representations
    of the kernel-operator structure.
\end{enumerate}

The result is a picture in which memory is not a property added to a
dynamics but the defining characteristic of dynamics as such. Systems
that remember their history are the rule; systems that forget are the
exception produced by setting $K = \delta$.

The formula sheet, revisited: it is not a list but a sample. Every
formula in it is a point in the continuous operator space parameterized
by kernel, order, and constraint structure. Fractional calculus is the
curve connecting two such points. RSVP, Yarncrawler, Spherepop, KES,
and Constraint Closure are structured regions of the same space.

The algebra of kernels is the algebra of memory. Memory is the algebra
of time.

%%=============================================================
\begin{thebibliography}{99}
%%=============================================================

\bibitem{oldham1974}
K.~B. Oldham and J.~Spanier,
\textit{The Fractional Calculus: Theory and Applications of
Differentiation and Integration to Arbitrary Order},
Academic Press, New York, 1974.

\bibitem{podlubny1999}
I.~Podlubny,
\textit{Fractional Differential Equations},
Academic Press, San Diego, 1999.

\bibitem{metzler2000}
R.~Metzler and J.~Klafter,
``The random walk's guide to anomalous diffusion: a fractional
dynamics approach,''
\textit{Physics Reports}, 339(1):1--77, 2000.

\bibitem{kilbas2006}
A.~A. Kilbas, H.~M. Srivastava, and J.~J. Trujillo,
\textit{Theory and Applications of Fractional Differential Equations},
North-Holland Mathematics Studies, vol.~204,
Elsevier, Amsterdam, 2006.

\bibitem{samko1993}
S.~G. Samko, A.~A. Kilbas, and O.~I. Marichev,
\textit{Fractional Integrals and Derivatives: Theory and Applications},
Gordon and Breach, Philadelphia, 1993.

\bibitem{mainardi2010}
F.~Mainardi,
\textit{Fractional Calculus and Waves in Linear Viscoelasticity},
Imperial College Press, London, 2010.

\bibitem{gorenflo2014}
R.~Gorenflo, A.~A. Kilbas, F.~Mainardi, and S.~V. Rogosin,
\textit{Mittag-Leffler Functions, Related Topics and Applications},
Springer, Berlin, 2014.

\bibitem{maclane1971}
S.~Mac~Lane,
\textit{Categories for the Working Mathematician},
Graduate Texts in Mathematics, vol.~5,
Springer, New York, 1971.

\bibitem{riehl2017}
E.~Riehl,
\textit{Category Theory in Context},
Dover Publications, Mineola, NY, 2017.

\bibitem{joyal1991}
A.~Joyal and R.~Street,
``The geometry of tensor calculus, I,''
\textit{Advances in Mathematics}, 88(1):55--112, 1991.

\bibitem{abramsky2004}
S.~Abramsky and B.~Coecke,
``A categorical semantics of quantum protocols,''
\textit{Proc.\ 19th IEEE Symposium on Logic in Computer Science},
415--425, 2004.

\bibitem{bredon1997}
G.~E. Bredon,
\textit{Sheaf Theory}, 2nd ed.,
Graduate Texts in Mathematics, vol.~170,
Springer, New York, 1997.

\bibitem{baez2010}
J.~C. Baez and M.~Stay,
``Physics, topology, logic and computation: a Rosetta Stone,''
in \textit{New Structures for Physics}, Lecture Notes in Physics
vol.~813, Springer, Berlin, 2010.

\bibitem{penrose2004}
R.~Penrose,
\textit{The Road to Reality},
Jonathan Cape, London, 2004.

\bibitem{flyxion-rsvp}
Flyxion,
``Adjunction and Irreversibility: RSVP Field Theory with Path Dependence,''
monograph, 129 pp., 2025.

\bibitem{flyxion-yarncrawler}
Flyxion,
``Yarncrawler: World-State Reconstruction via Sheaf-Variational
Equivalence and the Identifiability Theorem,''
monograph, 2025.

\bibitem{flyxion-constraint}
Flyxion,
``Constraint Before Completion: A Trajectory-First Theory of Work,''
essay, 2025.

\bibitem{flyxion-spherepop}
Flyxion,
``Spherepop: Categorical Semantics of an Irreversible Event Calculus,''
monograph (English and Spanish editions), 2025.

\bibitem{flyxion-compiled}
Flyxion,
``The Compiled Self: Identity as Fixed Point of Constraint-Preserving
Adjunction over Irreversible History Topos,'' essay, 2025.

\bibitem{flyxion-rsvp-positive}
Flyxion,
``RSVP and Positive Geometry: Deriving Associahedra and Amplituhedra
from Scalar-Vector-Energy Field Structure,'' v10, 38 pp., 2025.

\bibitem{flyxion-tartan}
Flyxion,
``TARTAN: Recursive Tiling, Ising Synchronization, and Emergent
Quantum Dynamics,'' monograph, 2025.

\bibitem{barandes2023}
J.~A. Barandes,
``The Unistochastic Reformulation of Quantum Theory,''
arXiv:2301.01609, 2023.

\bibitem{minkowski1908}
H.~Minkowski,
``Raum und Zeit,''
address delivered at the 80th Assembly of German Natural Scientists
and Physicians, Cologne, 1908. Translated in \textit{The Principle
of Relativity}, Dover, 1952.

\bibitem{misner1973}
C.~W. Misner, K.~S. Thorne, and J.~A. Wheeler,
\textit{Gravitation},
W.~H. Freeman, San Francisco, 1973.

\end{thebibliography}

\end{document}
